\providecommand{\cL}{\mathcal L}

\section{Bulk–boundary functoriality and spectral $R(z)$}
\label{sec:bulk-boundary-R}
\label{sec:spectral_braiding}
\label{sec:bulk-boundary}

\subsection{Boundary data and factorization module structure}
\label{subsec:boundary-module}
Let $A=(A_{\mathsf{ch}},A_{\mathsf{top}})$ be a $C_\ast\!\bigl(W(\mathsf{SC}^{\mathrm{ch,top}})\bigr)$–algebra.
A (topological) boundary condition along $t=0$ is modeled by a \emph{right $W(\mathsf{SC}^{\mathrm{ch,top}})$–module} $M$ supported on $\C\times \R_{\ge 0}$, i.e.\ a prefactorization algebra $\mathsf{Obs}^{\partial}$ on half‑rectangles with actions
\[
C_\ast\bigl(\FM_k(\C)\times E_1(m)\bigr)\otimes A_{\mathsf{ch}}^{\otimes k}\otimes M^{\otimes m}\longrightarrow M,
\]
compatible with $W$–coherence.  We refer to $M$ as a \emph{boundary factorization module}.

\subsection{Line operators and the spectral parameter}
\label{subsec:line-ops}
Boundary line operators $L$ are point‑supported defects along $t=0$ with insertion points $z\in\C$.
Given two parallel boundary lines $L_1(z_1)$ and $L_2(z_2)$ (with $z_1\neq z_2$), the bulk-to-boundary OPE in the strip determines a composition
\[
\mathsf{Comp}_{z_1,z_2} \colon L_1(z_1)\otimes L_2(z_2) \longrightarrow L_2(z_2)\otimes L_1(z_1),
\]
meromorphic in $z_{12}=z_1-z_2$ with logarithmic poles along $z_{12}=0$.
This composition is induced by the $W(\mathsf{SC}^{\mathrm{ch,top}})$–module structure and is computed by configuration integrals over $\FM_2(\C)$ with $E_1$–homotopies along the half‑interval.

\begin{definition}[Spectral braiding]
\label{def:spectral-braiding}
The \emph{spectral braiding} (or \emph{spectral $R$–matrix}) is the family of isomorphisms
\[
R_{L_1,L_2}(z)\colon L_1\otimes L_2 \xrightarrow{\ \sim\ } L_2\otimes L_1,
\qquad \text{for generic } z\in \C^\times \text{ (away from a discrete set of poles)},
\]
obtained from $\mathsf{Comp}_{z_1,z_2}$ by analytic continuation in $z=z_{12}$.
\end{definition}

\subsection{Yang–Baxter equation from $W$–coherence}
\label{subsec:YBE}
Consider three boundary lines $L_1,L_2,L_3$ at positions $z_1,z_2,z_3$.
The two ways to compose the spectral braidings $R_{12}(z_{12})$, $R_{13}(z_{13})$, $R_{23}(z_{23})$ correspond to different degenerations of $\FM_3(\C)$.

\begin{theorem}[Yang--Baxter equation; \ClaimStatusProvedHere]
\label{thm:YBE}
Let\/ $\cA$ be a logarithmic\/ $\SCchtop$-algebra.
The spectral braiding $R(z)$ satisfies the Yang–Baxter equation
\[
R_{12}(z_{12})\,R_{13}(z_{13})\,R_{23}(z_{23}) \;=\; R_{23}(z_{23})\,R_{13}(z_{13})\,R_{12}(z_{12})
\]
as morphisms $L_1\otimes L_2\otimes L_3 \to L_3\otimes L_2\otimes L_1$ for $z_i\neq z_j$.
\end{theorem}

\begin{proof}
We apply Stokes' theorem on $\FM_3(\C)$ to the logarithmic form $\omega_3^{\partial}$ governing the three-point bulk-to-boundary correlation of line operators $L_1(z_1), L_2(z_2), L_3(z_3)$.  By Definition~\ref{def:log-SC-algebra} and Theorem~\ref{thm:FM-calculus}, $\omega_3^{\partial}$ extends to a logarithmic form on $\FM_3(\C)$ satisfying $d\omega_3^{\partial} = 0$ in the interior $\Conf_3(\C)$.

\medskip
\noindent\textbf{Step 1: Boundary strata of $\FM_3(\C)$.}
By the FM boundary structure (Definition~\ref{def:boundary_divisors}), the codimension-1 faces of $\partial\FM_3(\C)$ are the collision divisors
\[
D_{\{1,2\}},\quad D_{\{1,3\}},\quad D_{\{2,3\}},\quad D_{\{1,2,3\}}.
\]
The full-collision stratum $D_{\{1,2,3\}} \cong \FM_3(\C)$ contributes a term proportional to the identity.  Here the module factorization data is already part of the $W(\SCchtop)$-module structure, and the same degree count shows that the restricted weight form is exact on $\FM_3(\C)$.  Thus only the three pairwise divisors contribute.

\medskip
\noindent\textbf{Step 2: Spectral parameter assignments on each face.}
Each divisor $D_{\{i,j\}} \cong \FM_2(\C) \times \FM_2(\C)$ parametrizes configurations where $z_i$ and $z_j$ collide while the third point stays separated.  The cluster center inherits the spectral parameter of the pair, and the residue factorizes the three-point correlator into a product of two-point braidings.  Explicitly:

\begin{itemize}
\item \textbf{Face $D_{\{1,2\}}$}: Points $z_1, z_2$ collide first.  The inner braiding carries spectral parameter $z_{12} = z_1 - z_2$ and produces $R_{12}(z_{12})$.  The cluster center (at $z_2$) then braids with $z_3$, producing $R_{13}(z_{13})$ followed by $R_{23}(z_{23})$.  The total contribution is
\[
D_{\{1,2\}}: \quad R_{12}(z_{12})\, R_{13}(z_{13})\, R_{23}(z_{23}).
\]

\item \textbf{Face $D_{\{2,3\}}$}: Points $z_2, z_3$ collide first, giving inner braiding $R_{23}(z_{23})$.  The remaining braidings with $z_1$ yield $R_{13}(z_{13})$ then $R_{12}(z_{12})$.  The total contribution is
\[
D_{\{2,3\}}: \quad R_{23}(z_{23})\, R_{13}(z_{13})\, R_{12}(z_{12}).
\]

\item \textbf{Face $D_{\{1,3\}}$}: Points $z_1, z_3$ collide.  This face contributes
\[
D_{\{1,3\}}: \quad R_{13}(z_{13})\, R_{12}(z_{12})\, R_{23}(z_{23}).
\]
\end{itemize}

\noindent\textbf{Step 3: Stokes' theorem and corner cancellation.}
By Stokes' theorem on the manifold with corners $\FM_3(\C)$,
\[
0 = \int_{\FM_3(\C)} d\omega_3^{\partial} = \sum_{\{i,j\}} (-1)^{\sigma_{\{i,j\}}} \int_{D_{\{i,j\}}} \Res_{D_{\{i,j\}}}(\omega_3^{\partial}).
\]
Codimension-2 corners arise from nested collisions (e.g., $D_{\{1,2\}} \cap D_{\{2,3\}}$, where $z_1 \to z_2$ and $z_2 \to z_3$ simultaneously).  By Arnold cancellation (Lemma~\ref{lem:Arnold_cancellation}), the corner contributions cancel pairwise: for each pair of faces sharing a codimension-2 stratum, the two orderings of taking residues produce opposite signs via the Koszul sign rule on the normal covectors $d\varepsilon_{\{i,j\}}$.

\medskip
\noindent\textbf{Step 4: Flatness and the Yang--Baxter equation.}
The Stokes identity $\int_{\partial\FM_3} = 0$ gives the
\emph{infinitesimal} three-point relation: the three pairwise
residues $r_{12}, r_{13}, r_{23}$ (at first order in $\hbar$) satisfy
$[r_{12}, r_{13}] + [r_{12}, r_{23}] + [r_{13}, r_{23}] = 0$,
which is the classical Yang--Baxter equation (CYBE).
All three faces $D_{\{1,2\}}, D_{\{1,3\}}, D_{\{2,3\}}$ contribute
to this three-term identity; none is discarded.

The \emph{quantum} Yang--Baxter equation
$R_{12}R_{13}R_{23} = R_{23}R_{13}R_{12}$ follows from the
flatness of the connection $\nabla$ on $\Conf_3(\C)$ whose
holonomy defines $R(z)$.  Flatness $\nabla^2 = 0$ is the content
of the Stokes argument above (the CYBE is the infinitesimal form
of $\nabla^2 = 0$).  A flat connection on $\Conf_3(\C)$ defines a
monodromy representation of the fundamental group
$\pi_1(\Conf_3(\C)) \cong PB_3$, the pure braid group on three
strands.  The pure braid group has the relation
$\sigma_{12}\sigma_{13}\sigma_{23} =
\sigma_{23}\sigma_{13}\sigma_{12}$ (this is the defining relation
of the pure braid group; cf.\ Birman \cite{Birman74}).  Applying
the monodromy representation $\rho\colon PB_3 \to
\mathrm{Aut}(L_1 \otimes L_2 \otimes L_3)$, with
$\rho(\sigma_{ij}) = R_{ij}(z_{ij})$, yields the quantum YBE.
The Stokes argument \emph{proves flatness}; the braid group
\emph{converts flatness to YBE}.
The path-ordered exponential defining $R(z)$ converges in the completed tensor product $\cA^! \widehat{\otimes}\, \cA^!((z^{-1}))$ because $r(z) \in z^{-1}\cA^! \widehat{\otimes}\, \cA^![[z^{-1}]]$ has no constant term: the Picard iteration $R^{(n)}(z) = 1 + \int_0^z r(\zeta) R^{(n-1)}(\zeta)\,d\zeta$ converges in the $z^{-1}$-adic topology.
\end{proof}

\begin{remark}[Yang--Baxter as associativity of Lagrangian convolution]
\label{rem:YBE-lagrangian-convolution}
The Yang--Baxter equation admits a geometric reading
through the Chriss--Ginzburg lens.
Let $\cL \hookrightarrow \cM$ be the Lagrangian inclusion
of the boundary into the ambient derived stack, and let
$\mathfrak{S}_b = \cL \times_{\cM} \cL$ be the Steinberg
self-intersection.
The $R$-matrix $R(z)$ is the monodromy of the flat connection
on $\mathfrak{S}_b$; the spectral parameter $z$ is the
coordinate on $\FM_2(\C)$---the closed-color factor of the
correspondence.
The three terms $R_{12}R_{13}R_{23}$ and $R_{23}R_{13}R_{12}$
correspond to the two ways of composing three Lagrangian
correspondences: they are the two iterated fiber products
\[
(\cL \times_{\cM} \cL) \times_{\cM} \cL
\;\simeq\;
\cL \times_{\cM} (\cL \times_{\cM} \cL).
\]
Associativity of the fiber product ensures the triple
self-intersection is well-defined; the \emph{equality}
$R_{12}R_{13}R_{23} = R_{23}R_{13}R_{12}$ requires the additional
input of \emph{flatness} of the spectral connection
(Theorem~\ref{thm:YBE}), which is the zero-curvature condition
underlying Theorem~\ref{thm:steinberg-presentation}(iii)(b).  The complete
strictification theorem
(Theorem~\ref{thm:complete-strictification}) proves this $R(z)$ is
a \emph{strict} factorization $R$-matrix: root multiplicity one
forces the spectral Drinfeld class to vanish at every filtration.
\end{remark}

\begin{proposition}[Lagrangian convolution interpretation of YBE;
\ClaimStatusProvedHere]
\label{prop:YBE-triple-lagrangian}
Let\/ $\cL \hookrightarrow \cM$ be the Lagrangian inclusion and\/
$\mathfrak{S}_b = \cL \times_{\cM}^h \cL$ the Steinberg
self-intersection.  The Yang--Baxter equation\/
$R_{12}R_{13}R_{23} = R_{23}R_{13}R_{12}$ \textup{(}proved
independently in Theorem~\textup{\ref{thm:YBE}} via Stokes on
$\FM_3(\C)$\textup{)} admits the following Lagrangian restatement:
\begin{enumerate}[label=\textup{(\roman*)}]
\item Line operators $L_i$ are modules for the convolution algebra
  $H_*^{\mathrm{BM}}(\mathfrak{S}_b)$ via the standard
  module identification: the line category is equivalent to
  $\mathfrak{S}_b$-modules
  \textup{(}Chriss--Ginzburg~\textup{\cite{CG97},} Chapter~8\textup{)}.
\item Through this identification, the spectral $R$-matrix
  $R_{ij}(z) \in \End(L_i \otimes L_j)$ corresponds to the
  monodromy of the spectral flat connection on $\mathfrak{S}_b$
  pulled back along the pair-projection
  $p_{ij}\colon \mathfrak{S}_b^{(3)} \to \mathfrak{S}_b$.
\item The YBE is the statement that the monodromy representation\/
  $\rho\colon PB_3 \to \Aut(L_1 \otimes L_2 \otimes L_3)$ respects
  the pure braid relation; this follows from flatness of the
  connection on $\Conf_3(\C)$.
\end{enumerate}
In particular, associativity of the fiber product
$(\cL \times_{\cM} \cL) \times_{\cM} \cL \simeq
\cL \times_{\cM} (\cL \times_{\cM} \cL)$
ensures the triple self-intersection $\mathfrak{S}_b^{(3)}$ is
well-defined, providing the geometric arena on which the flat
connection and its monodromy live.  The non-trivial content is
\emph{flatness}, not associativity.
\end{proposition}

\begin{proof}
Denote the triple self-intersection by
$\mathfrak{S}_b^{(3)} = \cL \times_{\cM} \cL \times_{\cM} \cL$.
This carries three projections to the pairwise self-intersections:
\[
  p_{ij} \colon \mathfrak{S}_b^{(3)} \longrightarrow
  \mathfrak{S}_b = \cL \times_{\cM} \cL,
  \qquad 1 \leq i < j \leq 3,
\]
given by projecting onto the $i$th and $j$th copies of $\cL$.
The spectral flat connection on $\mathfrak{S}_b$ (whose existence
follows from the OPE-Lagrangian correspondence of
Theorem~\ref{thm:bar-is-self-intersection}) pulls back via
$p_{ij}$ to connections on $\mathfrak{S}_b^{(3)}$.  These pullback
connections on the common space $\mathfrak{S}_b^{(3)}$ assemble
into a single flat connection $\nabla^{(3)}$ on $\Conf_3(\C)$
(the spectral parameter space).  Flatness of $\nabla^{(3)}$ is the
content of the Stokes argument in Theorem~\ref{thm:YBE}.

The monodromy representation of the flat connection $\nabla^{(3)}$
factors through $\pi_1(\Conf_3(\C)) \cong PB_3$.  Line operators
$L_i$ are $\mathfrak{S}_b$-modules (via the Chriss--Ginzburg
standard module: $H_*^{\mathrm{BM}}(\cL)$ is the standard module
for the convolution algebra $H_*^{\mathrm{BM}}(\mathfrak{S}_b)$,
and line operators are the spectral-parameter-dependent analog).
The $R$-matrix $R_{ij}(z_{ij})$ is the image of the pure braid
generator $\sigma_{ij}$ under the monodromy representation
$\rho\colon PB_3 \to \Aut(L_1 \otimes L_2 \otimes L_3)$.
Since $\rho$ is a group homomorphism and $PB_3$ satisfies the
relation $\sigma_{12}\sigma_{13}\sigma_{23} =
\sigma_{23}\sigma_{13}\sigma_{12}$, we conclude
$R_{12}R_{13}R_{23} = R_{23}R_{13}R_{12}$.
\end{proof}

\subsection{Classical limit and the infinitesimal $r(z)$}
\label{subsec:classical-r}
Let $\hbar$ denote the loop parameter of the perturbative HT theory.
Assume $R(z)$ admits an expansion
\[
R(z) \;=\; \mathrm{id} \,+\, \hbar\, r(z) \,+\, O(\hbar^2),
\]
in endomorphisms of $L_1\otimes L_2$.
Then $r(z)$ is a solution of the classical Yang–Baxter equation (CYBE) with spectral parameter:
\[
[ r_{12}(z_{12}), r_{13}(z_{13}) ] + [ r_{12}(z_{12}), r_{23}(z_{23}) ] + [ r_{13}(z_{13}), r_{23}(z_{23}) ] \;=\; 0.
\]

\begin{proposition}[Field--theoretic expression for $r(z)$;
\ClaimStatusProvedHere]
\label{prop:field-theory-r}
Let\/ $\cA$ be a logarithmic\/ $\SCchtop$-algebra.
Let $\{\cdot{}_\lambda \cdot\}$ be the bulk $(-1)$–shifted $\lambda$–bracket on cohomology (Section~\ref{sec:Ainfty-to-PVA}). The kernel of the infinitesimal spectral $r$–matrix is computed at first order by
\[
r(z) \;\simeq\; \int_{0}^{\infty} \! \mathrm{d}\lambda\ \mathrm{e}^{-\lambda z}\ \bigl\{\; \cdot {}_\lambda \cdot \;\bigr\},
\]
under the BV pairing identifying bulk and boundary modes, and with the propagator fixing the time‑ordering on the half‑line.  The integral converges for $\operatorname{Re}(z) > 0$ and extends meromorphically to $\C^\times$. In the formal algebraic setting, this identity holds as an equality of formal power series in $z^{-1}$.
\end{proposition}

\begin{proof}
We prove both the Laplace-transform formula and the classical Yang--Baxter equation in four steps.

\medskip
\noindent\textbf{Step 1: Perturbative expansion.}
Write $R(z) = \id + \hh\, r(z) + O(\hh^2)$ in $\End(L_1 \otimes L_2)$.  At tree level ($\hh^0$), line operators compose freely and $R = \id$.  The first quantum correction $r(z)$ is computed by a single bulk-to-boundary propagator exchange: a bulk operator propagates from position $z_1$ on line $L_1$ to position $z_2$ on line $L_2$ (with $z = z_1 - z_2$) through the half-space $\C \times \R_{\geq 0}$.

\medskip
\noindent\textbf{Step 2: Propagator integral.}
For physical realizations covered by Theorem~\ref{thm:physics-bridge}, the free propagator on $\C \times \R$ takes the form
\[
K(z,t) = \frac{\Theta(t)}{2\pi z}\, e^{-\mu(z) t} + (\text{regular}),
\]
where $\Theta(t)$ is the Heaviside function enforcing time-ordering and $\mu(z)$ encodes the holomorphic dependence.  The bulk $\lambda$-bracket kernel $K_\lambda$ is the Fourier-conjugate mode: $\{\cdot\,{}_\lambda\,\cdot\}$ acts on cohomology classes $[a], [b] \in H^\bullet(A,Q)$ via
\[
\{a_\lambda b\} = \int_0^\infty dt\; e^{\lambda t}\, K(z_{12}, t)\, \langle a \otimes b \rangle_{\mathrm{BV}},
\]
where $\langle \cdot \otimes \cdot \rangle_{\mathrm{BV}}$ is the BV pairing contracting bulk operators against boundary modes (cf.\ Section~\ref{sec:Ainfty-to-PVA}).

\medskip
\noindent\textbf{Step 3: Laplace transform.}
The first-order braiding $r(z)$ is obtained by evaluating the single-exchange diagram.  Performing the $t$-integration with the exponential weight $e^{-zt}$ from the spectral parameter gives
\[
r(z) = \int_0^\infty d\lambda\; e^{-\lambda z}\, \{\ \cdot{}_\lambda \cdot\ \}.
\]
This is the Laplace transform of the $\lambda$-bracket kernel, convergent for $\operatorname{Re}(z) > 0$ by the topological decay supplied in physical realizations by Theorem~\ref{thm:physics-bridge}.  Meromorphic continuation to $\C^\times$ follows from the meromorphy of the propagator in $z$.  In the abstract logarithmic $\SCchtop$ setting, the same identity is interpreted formally as a power series in $z^{-1}$.

For the free propagator $K(z,t) = \Theta(t)/(2\pi z)$, the $\lambda$-bracket is $\lambda$-independent and the Laplace integral evaluates to $r(z) \sim 1/z$ (consistent with the abelian CS result of Example~\ref{ex:Heisenberg_Yangian}).

\begin{remark}[Laplace transform as spectral duality]
\label{rem:laplace-spectral-duality}
The formula $r(z) = \int_0^\infty e^{-\lambda z} \{\cdot{}_\lambda \cdot\}\, d\lambda$ is the bridge between the PVA world (spectral variable~$\lambda$) and the $R$-matrix world (position variable~$z$).  Under DK-0, it recovers the classical $r$-matrix from the $\lambda$-bracket at the evaluation locus.  At the operadic level, this is the passage between the two colors of the Swiss-cheese operad: the $\lambda$-bracket encodes holomorphic (closed-color) singularities; $R(z)$ encodes topological (open-color) braiding.  The topological direction supplies the integration contour; the holomorphic direction supplies the meromorphic structure.  The Laplace transform \emph{is} the spectral duality of the HT theory.
\end{remark}

\begin{remark}[Specialization to the DK ladder]
\label{rem:spectral-R-dk-specialization}
On evaluation modules of simple $\fg$ at generic level $k$, the spectral $R$-matrix $R(z)$ equals the quantum group $R$-matrix of $U_q(\fg)$ at $q = e^{i\pi/(k+h^\vee)}$.  This is Theorem~\ref{thm:affine-monodromy-identification}: one-loop exactness of the affine half-space BV theory collapses the $A_\infty$ tower, the reduced HT connection is the KZ connection (Theorem~\ref{thm:reduced-equals-kz}), and the Drinfeld--Kohno theorem~\cite{Dri89, KL93} identifies the monodromy.  The reduced HT line category on evaluation modules is $\mathrm{Rep}_q(\fg)$ (Corollary~\ref{cor:affine-line-category}).
\end{remark}

\medskip
\noindent\textbf{Step 4: Classical Yang--Baxter equation and antisymmetry.}
Substituting $R(z) = \id + \hh\, r(z) + O(\hh^2)$ into the quantum YBE (Theorem~\ref{thm:YBE}), expanding to order $\hh^2$, and using $R_{ij} R_{ik} R_{jk} = R_{jk} R_{ik} R_{ij}$:
\begin{align*}
&(\id + \hh r_{12} + O(\hh^2))(\id + \hh r_{13} + O(\hh^2))(\id + \hh r_{23} + O(\hh^2)) \\
&\quad = (\id + \hh r_{23} + O(\hh^2))(\id + \hh r_{13} + O(\hh^2))(\id + \hh r_{12} + O(\hh^2)).
\end{align*}
At order $\hh^1$, both sides give $\id + \hh(r_{12} + r_{13} + r_{23})$, which is consistent.  At order $\hh^2$, the LHS contains $r_{12}r_{13} + r_{12}r_{23} + r_{13}r_{23}$ and the RHS contains $r_{23}r_{13} + r_{23}r_{12} + r_{13}r_{12}$.  Their difference yields
\[
[r_{12}(z_{12}),\, r_{13}(z_{13})] + [r_{12}(z_{12}),\, r_{23}(z_{23})] + [r_{13}(z_{13}),\, r_{23}(z_{23})] = 0,
\]
which is the classical Yang--Baxter equation with spectral parameter.

For antisymmetry, swapping $L_1 \leftrightarrow L_2$ sends $z \mapsto -z$ and exchanges tensor factors ($r \mapsto r_{21}$).  The spectral braiding satisfies $R_{12}(z) R_{21}(-z) = \id$ (inverse relation from the definition of braiding), so at order $\hh$:
\[
r_{12}(z) + r_{21}(-z) = 0, \qquad\text{i.e.,}\quad r(-z) = -r_{21}(z). \qedhere
\]
\end{proof}

\subsection{Koszul dual Hopf picture and factorisation quantum groups}
\label{subsec:Hopf}
Let $\mathbf B$ be the bar cooperad of $\mathsf{SC}^{\mathrm{ch,top}}$ and let $A$ be a $W(\mathsf{SC}^{\mathrm{ch,top}})$–algebra.
The bar construction $\mathsf{Bar}^{\mathrm{ch,top}}(A)$ is a colored coalgebra over $\mathbf B$; its cobar $\mathsf{Cobar}^{\mathrm{ch,top}}$ yields a (filtered) Hopf‑like object $\mathcal H$ encoding the OPE and higher products.

\begin{theorem}[Braiding on boundary categories; \ClaimStatusProvedHere]
\label{thm:braided-category}
Let\/ $\cA$ be a logarithmic\/ $\SCchtop$-algebra.
Let $\mathcal C_{\partial}$ be the (derived) category of boundary line operators with the tensor structure induced by the $W$–module structure of~$M$. Then $R(z)$ defines a meromorphic braiding on $\mathcal C_{\partial}$, and the classical limit $r(z)$ integrates to a filtered quasi‑triangular structure on the Koszul–dual Hopf object $\mathcal H$ (the "factorisation quantum group" in this setting).
\end{theorem}

\begin{proof}
We establish the meromorphic braiding and the quasi-triangular Hopf structure separately.

\medskip
\noindent\textbf{Part A: Meromorphic braiding on $\mathcal{C}_\partial$.}
The derived category $\mathcal{C}_\partial$ of boundary line operators inherits a monoidal structure from the $E_1$-composition along the boundary half-line $\R_{\geq 0}$: given right $W(\SCchtop)$-modules $M_1, M_2$ (boundary factorization modules in the sense of Section~\ref{subsec:boundary-module}), their tensor product $M_1 \otimes_{E_1} M_2$ is defined by the operadic composition in the open color.  The spectral braiding $R(z)$ of Definition~\ref{def:spectral-braiding} provides a natural isomorphism
\[
\beta_z \colon M_1 \otimes_{E_1} M_2 \xrightarrow{\ \sim\ } M_2 \otimes_{E_1} M_1
\]
for generic $z \in \C^\times$, arising from the exchange of insertion points in $\FM_2(\C)$.  By Theorem~\ref{thm:YBE}, $\beta_z$ satisfies the braid relation on triple tensor products, making $(\mathcal{C}_\partial, \otimes_{E_1}, \beta_z)$ a braided monoidal category with meromorphic dependence on $z$.

\medskip
\noindent\textbf{Part B: Filtered quasi-triangular structure on $\mathcal{H}$.}
Let $\mathcal{H} = \mathsf{Cobar}^{\mathrm{ch,top}}(\mathsf{Bar}^{\mathrm{ch,top}}(A))$ be the Hopf-like object from the bar--cobar construction.  By Theorem~\ref{thm:filtered-koszul}, $\mathcal{H}$ carries the holomorphic weight filtration $F^\bullet$ with $\gr^F \mathcal{H} \cong U(\mathfrak{g}[z])$ (the enveloping algebra of the loop algebra, on the closed color) tensored with an associative algebra (on the open color).

The classical $r$-matrix $r(z)$ defines an element $r(z) \in \mathcal{H} \widehat{\otimes}\, \mathcal{H}((z^{-1}))$ satisfying:
\begin{enumerate}[label=(\alph*)]
\item The CYBE (Proposition~\ref{prop:field-theory-r}, Step 4 of the proof), which is the infinitesimal form of the braid relation;
\item Antisymmetry $r(-z) = -r_{21}(z)$ (Theorem~\ref{thm:spectral_R_YBE}(iii));
\item The intertwining property $r(z)\,\Delta(x) - \Delta^{\mathrm{op}}(x)\,r(z) = 0$ at leading order in $\hh$ (from the quasi-triangularity established in the proof of Theorem~\ref{thm:Koszul_dual_Yangian}).
\end{enumerate}
These three properties define a \emph{filtered quasi-triangular structure}: $r(z)$ lives in filtration degree $F^1(\mathcal{H} \otimes \mathcal{H})$, and the CYBE holds modulo $F^2$.  The full quantum $R$-matrix $R(z) = \id + \hh\, r(z) + O(\hh^2)$ extends this to a quasi-triangular structure on the completed object $\widehat{\mathcal{H}}$, with the Yang--Baxter equation holding to all orders by Theorem~\ref{thm:YBE}.

The bar--cobar duality (Theorem~\ref{thm:bar-cobar-adjunction}) ensures that representations of $\mathcal{H}$ are equivalent to $\SCchtop$-modules, so the braided monoidal category $\mathcal{C}_\partial$ from Part A is identified with the category of $\mathcal{H}$-modules equipped with the braiding induced by $R(z)$.
\end{proof}

\subsection{OPE of line operators and the spectral $R$-matrix}

\subsubsection{Meromorphic tensor product}

Two line operators $\ell_1(z_1)$, $\ell_2(z_2)$ at distinct holomorphic positions compose via a meromorphic tensor product:
\begin{equation}
\ell_1(z_1) \ell_2(z_2) \sim \sum_k C_k(z_1 - z_2) \, \ell_k,
\end{equation}
where $C_k(z)$ are meromorphic functions in $z$ encoding quantum corrections.

\begin{definition}[Spectral $R$-Matrix]
\label{def:spectral_R}
The \emph{spectral $R$-matrix} is a degree-1 element:
\begin{equation}
r(z) \in \mathcal{A}^! \widehat{\otimes} \mathcal{A}^!((z^{-1})),
\end{equation}
encoding the quantum corrections to the tensor product of line operators:
\begin{equation}
\ell_1(z) \otimes \ell_2(0) \sim e^{r(z)} (\ell_1 \otimes \ell_2).
\end{equation}
\end{definition}

\begin{remark}[Three guises of $r(z)$]
\label{rem:r-z-disambiguation}
The symbol $r(z)$ is used in three equivalent guises: (a) the degree-$1$ element $r(z) \in \cA^! \widehat{\otimes}\, \cA^!((z^{-1}))$ encoding quantum corrections to the line-operator tensor product (Definition~\ref{def:spectral_R}); (b) the arity-$(2,0)$ projection of the MC element $\alpha_T$ (the closed-color binary datum); (c) the Laplace transform $r(z) = \int_0^\infty e^{-\lambda z}\{{\cdot}\,{}_\lambda\,{\cdot}\}\,d\lambda$ of the PVA $\lambda$-bracket (Proposition~\ref{prop:field-theory-r}). These agree by Theorem~\ref{thm:resolvent-principle}(iii).
\end{remark}

\begin{remark}[Relation between $r(z)$ and $R(z)$]
\label{rem:r-vs-R-relation}
For Koszul algebras ($m_k^H = 0$ for $k \geq 3$,
Corollary~\ref{cor:koszul-resolvent}), the braiding
isomorphism is $R(z) = \exp(r(z))$ in the completed tensor
product.  In the general (non-formal Swiss-cheese) case, the higher
$\Ainf$ operations $m_k^H$ ($k \geq 3$) contribute
corrections to the spectral connection form
(Proposition~\ref{prop:spectral-resolvent}), and the
$R$-matrix is the path-ordered exponential (holonomy) of this
corrected connection rather than the ordinary exponential of a
single element.
\end{remark}

\subsubsection{Properties of $r(z)$}

\begin{theorem}[Spectral $R$-Matrix Satisfies $A_\infty$ Yang--Baxter;
\ClaimStatusProvedElsewhere]
\label{thm:spectral_R_YBE}
Let\/ $\cA$ be a logarithmic\/ $\SCchtop$-algebra.
(Dimofte--Niu--Py \cite{DNP25}, Theorem 5.2) The spectral $R$-matrix $r(z)$ satisfies:
\begin{enumerate}[label=(\roman*)]
\item \textbf{Maurer--Cartan equation}:
\begin{equation}
Q(r(z)) + \frac{1}{2}[r(z), r(z)] + \sum_{k \geq 3} m_k(r(z)^{\otimes k}) = 0;
\end{equation}

\item \textbf{$A_\infty$ Yang--Baxter equation}:
\begin{equation}
[r_{12}(w), r_{13}(z+w)] + [r_{12}(w), r_{23}(z)] + [r_{13}(z+w), r_{23}(z)] + (\text{higher $m_k$ terms}) = 0,
\end{equation}
in $\mathcal{A}^{\otimes 3}[[z^{-1}, w^{-1}, (z-w)^{-1}]]$;

\item \textbf{Skew-symmetry}:
\begin{equation}
r(-z) = -r_{21}(z),
\end{equation}
where $r_{21}$ means swapping the two tensor factors.
\end{enumerate}
\end{theorem}

\begin{proof}[Proof sketch]
\textbf{(i) Maurer--Cartan}: This follows from the associativity of the OPE. The OPE of three lines $\ell_1(z_1) \ell_2(z_2) \ell_3(z_3)$ must be independent of the order of expansion. This forces $r(z)$ to satisfy the MC equation in $\mathcal{A}^! \otimes \mathcal{A}^!$.

\textbf{(ii) $A_\infty$ Yang--Baxter}: The associativity of the OPE for three lines yields:
\begin{equation}
r_{23}(z) + (id \otimes \Delta_z)(r(z+w)) = r_{12}(w) + (\Delta_w \otimes id)(r(z)),
\end{equation}
where $\Delta_z : \mathcal{A}^! \to \mathcal{A}^! \otimes_{r(z)} \mathcal{A}^!$ is the coproduct deformed by $r(z)$. Expanding this identity and using the $A_\infty$ structure yields the $A_\infty$ Yang--Baxter equation.

\textbf{(iii) Skew-symmetry}: This follows from the symmetry of the tensor product: swapping $\ell_1 \leftrightarrow \ell_2$ corresponds to $z \to -z$ and exchanging tensor factors.
\end{proof}

\subsection{dg-Shifted Yangians from operadic Koszul duality}
\label{subsec:dg-yangian-operadic}

The dg-shifted Yangian axioms are not posited \emph{ad hoc}.
They are the operadic identities of an
$\SCchtop$-algebra, decomposed by color.  The link is the
Koszul self-duality of~$\SCchtop$: the operadic Koszul dual
of the chiral Swiss-cheese operad is itself, so
the Koszul dual~$\cA^!$ of an $\SCchtop$-algebra~$\cA$ is
\emph{again} an $\SCchtop$-algebra.  The Yangian axioms are
the explicit presentation of this inherited structure.

The convolution algebra on the Steinberg variety
\emph{is} the Hecke algebra, by the structure of the
correspondence.  Likewise, the $\SCchtop$-algebra structure
on~$\cA^!$ \emph{is} the dg-shifted Yangian, by the
structure of the two-colored operad.

\subsubsection*{Koszul self-duality of the Swiss-cheese operad}

\begin{theorem}[Koszul self-duality of $\SCchtop$;
\ClaimStatusProvedHere]
\label{thm:SC-self-duality}
\index{Swiss-cheese operad!Koszul self-duality|textbf}
\index{Koszul duality!Swiss-cheese self-duality}
The operadic Koszul dual of\/ $\SCchtop$ is quasi-isomorphic
to\/ $\SCchtop$ itself:
\[
  (\SCchtop)^! \;\simeq\; \SCchtop
\]
as two-colored dg operads \textup(up to operadic
desuspension\textup).  In particular, for any
$\SCchtop$-algebra~$\cA$, the linear dual\/
$\cA^! := \Hom(\barB^{\SCchtop}(\cA),\, k)$
is again an\/ $\SCchtop$-algebra.
\end{theorem}

\begin{proof}
The proof has three steps: component self-duality,
interaction self-duality, and transfer via formality.

\medskip
\noindent\textbf{Step 1: Component self-duality.}
The closed-color operad~$E_2$ is Koszul with
$E_2^! \simeq E_2$ (Ginzburg--Kapranov~\cite{GK94}):
the Koszul dual of the commutative operad is the Lie operad,
and both sit inside~$E_2$ in a self-dual configuration.
The open-color operad $\mathsf{Ass} = E_1$ is Koszul
with $\mathsf{Ass}^! \simeq \mathsf{Ass}$
(Ginzburg--Kapranov~\cite{GK94}): the associative operad is
self-dual.  Both self-dualities hold up to the standard
operadic desuspension~$\mathcal{S}^{-1}$.

\medskip
\noindent\textbf{Step 2: Interaction self-duality.}
The classical Swiss-cheese operad~$\mathsf{SC}$ is built
from~$E_2$ and~$\mathsf{Ass}$ by a distributive law
$\lambda \colon E_2 \circ \mathsf{Ass} \to
\mathsf{Ass} \circ E_2$
governing the mixed operations
(Voronov~\cite{Vor99}).  By Livernet~\cite{Liv06},
Proposition~5.4, $\mathsf{SC}$ is Koszul and
$\mathsf{SC}^!$ is built from $E_2^!$ and
$\mathsf{Ass}^!$ by the dual distributive
law~$\lambda^!$.  The key input: the
quadratic relation encoding ``the closed-color bracket
is a derivation of the open-color product'' is
self-orthogonal under the operadic pairing.  Since the
component operads are self-dual and the distributive law
dualizes to itself, the full Swiss-cheese operad is
self-dual: $\mathsf{SC}^! \simeq \mathsf{SC}$.

Concretely, the mixed quadratic relation is the Leibniz
rule:
\begin{equation}\label{eq:leibniz-self-dual}
\mu(\beta(a,b),\, c)
\;=\;
\beta(\mu(a,c),\, b) + \beta(a,\, \mu(b,c)),
\end{equation}
where $\mu$ is the open-color product and $\beta$ is the
closed-color bracket.  Under the orthogonal complement
operation on mixed quadratic relations (which pairs
$(\mathsf{ch},\mathsf{ch},\mathsf{top};\mathsf{top})$
against $(\mathsf{ch},\mathsf{ch},\mathsf{top};
\mathsf{top})$ via the operadic inner product on two-colored
binary operations), the Leibniz
rule~\eqref{eq:leibniz-self-dual}
spans a subspace equal to its own orthogonal
complement---the derivation condition is the unique mixed
quadratic relation compatible with both $E_2$ and
$\mathsf{Ass}$ (Livernet~\cite{Liv06}, Proposition~5.4
(self-duality of SC); cf.~Vallette~\cite{Val07}, \S4.2).

\medskip
\noindent\textbf{Step 3: Transfer via formality.}
The formality quasi-isomorphism
$\Phi \colon \SCchtop \xrightarrow{\;\sim\;}
\mathsf{SC}$ (Theorem~\ref{thm:homotopy-Koszul}, Step~2)
is a map of two-colored dg operads.  The bar and cobar
functors for colored dg operads preserve
quasi-isomorphisms between levelwise cofibrant operads
over a field of characteristic zero (\cite{BM03,GK94}).
Composing:
\begin{equation}\label{eq:self-duality-chain}
(\SCchtop)^!
\;\xleftarrow[\;\sim\;]{\Omega\mathbf{B}(\Phi)^{-1}}\;
\mathsf{SC}^!
\;\xrightarrow[\;\sim\;]{\text{Step 2}}\;
\mathsf{SC}
\;\xleftarrow[\;\sim\;]{\Phi}\;
\SCchtop.
\end{equation}
The composite is the self-duality quasi-isomorphism.  The
arrows are compatible with the color decomposition: on the
closed color, the chain reduces to the Kontsevich formality
$C_\ast(\FM_\bullet(\C)) \simeq E_2 \simeq E_2^!$; on the
open color, it is the identity
$E_1 = \mathsf{Ass} = \mathsf{Ass}^!$; on the mixed
sector, it is the transfer of the self-dual distributive
law through formality.
\end{proof}

\begin{remark}[The Steinberg parallel]
\label{rem:self-duality-steinberg}
\index{Steinberg variety!operadic self-duality}
The self-duality $(\SCchtop)^! \simeq \SCchtop$ reflects
the symmetry of the bar complex on
$\FM_k(\C) \times \Conf_k(\R)$ viewed as a two-sided
correspondence: convolution (differential) and
diagonal (coproduct) are defined by the same
operadic structure, so dualizing the
correspondence reproduces the same type of algebra.
\end{remark}

\subsubsection*{Recognition: the Yangian axioms as operadic
identities}

\begin{theorem}[Yangian recognition;
\ClaimStatusProvedHere]
\label{thm:yangian-recognition}
\index{dg-shifted Yangian!recognition theorem|textbf}
\index{Swiss-cheese operad!Yangian recognition}
Let\/ $\mathcal{Y}$ be an\/ $\SCchtop$-algebra concentrated on a
single object \textup(i.e.\
$\mathcal{Y}_{\mathsf{ch}} = \mathcal{Y}_{\mathsf{top}} = \mathcal{Y}$ with
structure maps
$C_\ast(\FM_k(\C)) \otimes \mathcal{Y}^{\otimes k}
\to \mathcal{Y}$ and\
$E_1(m) \otimes \mathcal{Y}^{\otimes m}
\to \mathcal{Y}$
satisfying the\/ $\SCchtop$ operadic identities\textup).
Then\/ $\mathcal{Y}$ admits a canonical dg-shifted Yangian structure
$(\tau_z,\, r(z),\, \Delta_z,\, \varepsilon)$ in the sense
of Definition~\textup{\ref{def:dg_Yangian}}, and conversely
every dg-shifted Yangian arises in this way.
\end{theorem}

\begin{proof}
We construct a bijective dictionary between the operadic
data and the Yangian data.  The dictionary has four entries
(one for each piece of Yangian structure) and four
verifications (matching the Yangian axioms to operadic
identities).

\medskip
\noindent\textbf{Translation from the $E_2$ action
(closed color).}
The group $(\C, +)$ acts on $\FM_k(\C)$ by simultaneous
translation $z_i \mapsto z_i + w$.  This action is
encoded in the closed-color operadic structure: the induced
derivation on $\mathcal{Y}$ is the translation generator
$T = \partial_z$, and $\tau_z = \exp(zT)$ is the
$\C$-equivariance of the closed-color structure maps.
The degree of~$T$ is $(0,\text{odd},1)$: the
holomorphic-weight grading assigns weight~$1$ to
$\partial_z$.

\medskip
\noindent\textbf{Spectral $R$-matrix from $\FM_2(\C)$
(closed color, arity $2$).}
The closed-color structure map at arity~$2$ is
$C_\ast(\FM_2(\C)) \otimes \mathcal{Y}^{\otimes 2} \to \mathcal{Y}$.
The Fulton--MacPherson space
$\FM_2(\C) \simeq S^1$ is parametrized by
the relative coordinate $z = z_1 - z_2$; the fundamental
class $[\FM_2] \in C_1(\FM_2(\C))$ and its boundary
strata encode the spectral $R$-matrix
$r(z) \in \mathcal{Y} \widehat{\otimes}\, \mathcal{Y}((z^{-1}))$
of degree $(1,\text{odd},0)$.  The homological degree is~$1$
from $\dim_\R \FM_2(\C) = 1$ (equivalently, from the bar
desuspension); the odd parity is the fermion number; the
holomorphic weight is~$0$ because the leading pole $1/z$ is
the fundamental class contribution.

\medskip
\noindent\textbf{Coproduct and counit from $E_1$
(open color).}
The open-color $E_1$-algebra structure gives an
associative product on~$\mathcal{Y}$.  Passing to the
dual coalgebra description: the cofree coalgebra
$T^c(s^{-1}\bar{\mathcal{Y}})$ carries the deconcatenation
coproduct (Proposition~\ref{prop:bar-correspondence}(b));
dualizing produces the twisted coproduct
$\Delta_z \colon \mathcal{Y} \to
\mathcal{Y} \otimes_{r(z)} \mathcal{Y}[[z^{-1},z]]$
on $\mathcal{Y}$ (the $z$-dependence enters from the interaction
with the closed color; at the $E_0$-page of the holomorphic
weight filtration, $\Delta_z$ reduces to the $z$-independent
deconcatenation).  The counit
$\varepsilon \colon \mathcal{Y} \to \C$ is the augmentation of the
$E_1$-algebra: evaluation on the vacuum module.

\medskip
\noindent\textbf{Yangian axioms from operadic identities.}
The operadic identities for $\SCchtop$ at mixed color
profiles decompose, entry by entry, into the axioms of
Definition~\ref{def:dg_Yangian}:
\begin{enumerate}[label=\textup{(\alph*)}]
\item \emph{Counit axioms.}
  $(\varepsilon \otimes \id) \circ \Delta_z = \id$
  and $(\id \otimes \varepsilon) \circ \Delta_z
  \simeq \tau_z$
  are the operadic unit identities for the open color,
  composed with the $\C$-equivariance of the closed color:
  removing one line recovers the other, shifted by the
  holomorphic separation~$z$.
\item \emph{Translation invariance.}
  $(\tau_w \otimes \tau_w)\, r(z) = r(z)$ is the statement
  that $r(z)$ depends only on $z = z_1 - z_2$, not on
  individual positions---the $\C$-equivariance of
  $\FM_2(\C)$.
  The identity $\varepsilon(T(a)) = 0$ is the
  $\C$-invariance of the vacuum.
\item \emph{Weak co-commutativity.}
  $r(-z) = -r_{21}(z)$ is the involution
  $\sigma \colon \FM_2(\C) \to \FM_2(\C)$ swapping the two
  points ($z \mapsto -z$), composed with the swap
  $\mathcal{Y} \otimes \mathcal{Y} \to \mathcal{Y} \otimes \mathcal{Y}$.
  This is the $\Z/2$-equivariance of the arity-$2$
  closed-color operation.
\item \emph{Quantum $\Ainf$ Yang--Baxter.}
  The identity
  \[
    r_{23}(z) + (\id \otimes \Delta_z)(r(z{+}w))
    = r_{12}(w) + (\Delta_w \otimes \id)(r(z))
  \]
  is the operadic associativity constraint at arity~$3$: the
  two ways to compose the closed-color arity-$2$ operation
  (contributing~$r$) with the open-color arity-$2$ operation
  (contributing~$\Delta$) must agree by the
  $\SCchtop$ operadic identity at the mixed
  profile
  $(\mathsf{ch},\mathsf{ch},\mathsf{ch};\mathsf{top})$
  restricted to
  $\FM_3(\C) \times E_1(3)$.  This is the content of the
  $A_\infty$ Yang--Baxter equation
  (Corollary~\ref{cor:ainfty-ybe-as-convergence}).
\end{enumerate}

\medskip
\noindent\textbf{Converse.}
Given a dg-shifted Yangian
$(\mathcal{Y},\, \tau_z,\, r(z),\, \Delta_z,\, \varepsilon)$,
define $\SCchtop$-algebra structure maps by reversing the
dictionary: the closed-color maps at arity~$k$ are
determined by the $k$-fold iterated $r$-matrix (successive
pairwise OPEs parametrized by $\FM_k(\C)$); the open-color
maps are the $E_1$-product from the coproduct; the mixed
maps are the compatibility data.  The Yangian axioms
(Definition~\ref{def:dg_Yangian}) ensure that the operadic
identities are satisfied.  The higher-arity identities follow
because $\SCchtop$ is generated as a colored operad by binary
operations (arity-$2$ closed-color and arity-$2$ open-color):
every $k$-ary operadic identity decomposes into iterated
compositions of binary operations, and these compositions are
determined by the arity-$2$ and arity-$3$ data verified above
(Loday--Vallette~\cite{LV12}, Proposition~5.3.2).
\end{proof}

Having derived the Yangian axioms from the operadic
structure, we record the explicit form for reference.

\begin{definition}[dg-Shifted Yangian]
\label{def:dg_Yangian}
A \emph{dg-shifted Yangian} is an $A_\infty$ algebra
$\mathcal{Y}$ equipped with:
\begin{enumerate}[label=(\roman*)]
\item A translation automorphism
  $\tau_z : \mathcal{Y} \to \mathcal{Y}[[z]]$, generated by
  a derivation $T$ of degree $(0, \text{odd}, 1)$:
\begin{equation}
\tau_z = \exp(z T) = \sum_{n \geq 0} \frac{z^n}{n!} T^n;
\end{equation}

\item A spectral $R$-matrix
  $r(z) \in \mathcal{Y} \widehat{\otimes}
  \mathcal{Y}[[z^{-1}, z]]$ of degree
  $(1, \text{odd}, 0)$;

\item A twisted coproduct
  $\Delta_z : \mathcal{Y} \to \mathcal{Y} \otimes_{r(z)}
  \mathcal{Y}[[z^{-1}, z]]$ (an $A_\infty$ algebra
  morphism);

\item A counit $\varepsilon : \mathcal{Y} \to \C$ (also an
  $A_\infty$ morphism);
\end{enumerate}
satisfying:
\begin{itemize}
\item Counit axioms: $\varepsilon(1) = 1$,
  $(\varepsilon \otimes \id) \circ \Delta_z = \id$,
  $(\id \otimes \varepsilon) \circ \Delta_z
  \simeq \tau_z$;
\item Translation invariance: $\varepsilon(T(a)) = 0$,
  $\varepsilon(\tau_z(a)) = \varepsilon(a)$ for all
  $a \in \mathcal{Y}$;
\item Weak co-commutativity:
  $(\tau_w \otimes \tau_w) r(z) = r(z)$,
  $r(-z) = -r_{21}(z)$;
\item Co-associativity (quantum $A_\infty$ Yang--Baxter):
\begin{equation}
r_{23}(z) + (\id \otimes \Delta_z)(r(z+w))
= r_{12}(w) + (\Delta_w \otimes \id)(r(z)).
\end{equation}
\end{itemize}
By Theorem~\textup{\ref{thm:yangian-recognition}}, a
dg-shifted Yangian is precisely an\/ $\SCchtop$-algebra
concentrated on a single object.
\end{definition}

The Yangian identification of $\cA^!$ is now a structural
consequence, not a verification:

\begin{theorem}[Koszul dual is a dg-shifted Yangian;
\ClaimStatusProvedHere]
\label{thm:Koszul_dual_Yangian}
\index{Koszul dual!dg-shifted Yangian|textbf}
\index{dg-shifted Yangian!from Koszul duality}
For a chirally Koszul logarithmic\/ $\SCchtop$-algebra~$\cA$
\textup{(}Definition~\textup{\ref{def:log-SC-algebra})},
the Koszul dual~$\cA^!$ is a dg-shifted Yangian.
\end{theorem}

\begin{proof}
\textbf{Step 1: $\barB(\cA)$ is an
$\SCchtop$-coalgebra.}
The bar complex $\barB^{\SCchtop}(\cA)
= T^c(s^{-1}\bar{\cA})$ carries the bar differential
$d_{\barB}$ (a coderivation, from the closed color
$\FM_k(\C)$) and the deconcatenation coproduct~$\Delta$
(from the open color $\Conf_k(\R)$).  These assemble into
an $\SCchtop$-coalgebra structure by
Proposition~\ref{prop:bar-correspondence}: the differential
is the convolution and the coproduct is the diagonal of
the two-sided correspondence.

\medskip
\noindent\textbf{Step 2: Operadic dualization.}
By the Koszul self-duality of $\SCchtop$
(Theorem~\ref{thm:SC-self-duality}), the linear dual
$\cA^! = \Hom(\barB^{\SCchtop}(\cA),\, k)$ inherits an
$\SCchtop$-algebra structure.  The dictionary between
coalgebra data and algebra data is:
\begin{center}
\renewcommand{\arraystretch}{1.15}
\begin{tabular}{lll}
\textbf{Bar coalgebra} & $\longleftrightarrow$
  & \textbf{Koszul dual algebra} \\
\hline
Coproduct $\Delta$ (open color)
  && Product on $\cA^!$ (Yangian product) \\
Differential $d_{\barB}$ (closed color)
  && $\lambda$-bracket on $\cA^!$ (Sklyanin) \\
Coderivation property
  && $\lambda$-bracket is a derivation
     of the product
\end{tabular}
\end{center}

\medskip
\noindent\textbf{Step 3: Recognition.}
By Theorem~\ref{thm:yangian-recognition}, this
$\SCchtop$-algebra structure on $\cA^!$ is precisely a
dg-shifted Yangian.  The full identification:
\begin{center}
\renewcommand{\arraystretch}{1.15}
\begin{tabular}{lll}
\textbf{Operadic datum} & \textbf{Yangian datum}
  & \textbf{Origin} \\
\hline
$\C$-equivariance on $\FM_k(\C)$ &
  $\tau_z = \exp(zT)$ &
  holomorphic translation \\
$[\FM_2(\C)]$ at arity~$2$ &
  $r(z) \in \cA^! \widehat\otimes\, \cA^!((z^{-1}))$ &
  OPE of parallel lines \\
$E_1$-coalgebra structure &
  $\Delta_z,\; \varepsilon$ &
  ordered interval-cutting \\
Mixed operadic identity at $\FM_3 \times E_1(3)$ &
  quantum $\Ainf$ YBE &
  three-line consistency
\end{tabular}
\end{center}
\end{proof}

\begin{remark}[Logical status]
\label{rem:yangian-logical-status}
\index{dg-shifted Yangian!logical status}
The argument of Dimofte--Niu--Py~\cite{DNP25}
(Theorem~5.5) verified the Yangian axioms by direct
physical construction: translation from holomorphic
translation, $R$-matrix from the OPE, coproduct from
parallel lines.  The structural proof above replaces
verification with derivation: the Yangian axioms \emph{are}
the $\SCchtop$ operadic identities
(Theorem~\ref{thm:yangian-recognition}), so they hold by the
operadic bar-cobar duality that produces~$\cA^!$ in the
first place.  The physical constructions of \cite{DNP25}
are now \emph{identifications}---naming the operadic data in
physical language---not independent proofs.
\end{remark}

\subsubsection*{The duality involution}

\begin{theorem}[Duality involution;
\ClaimStatusProvedHere]
\label{thm:duality-involution}
\index{Koszul duality!involution|textbf}
\index{dg-shifted Yangian!duality involution}
Let\/ $\cA$ be a chirally Koszul logarithmic\/
$\SCchtop$-algebra.  Then
$(\cA^!)^! \simeq \cA$ as\/ $\SCchtop$-algebras.
The duality involution decomposes by color:
\begin{enumerate}[label=\textup{(\roman*)}]
\item \emph{Closed-color involution.}
  The chiral Koszul dual of $\cA^!$ \textup(with its
  Sklyanin $\lambda$-bracket\textup) recovers $\cA$'s
  chiral algebra structure:
  $\Omegach(\barBch(\cA^!)) \simeq \cA$.
\item \emph{Open-color involution.}
  The $E_1$-Koszul dual of $\cA^!$ \textup(with its
  Yangian product\textup) recovers $\cA$'s $E_1$-algebra
  structure:
  $\Omega^{\mathrm{ord}}(\barB^{\mathrm{ord}}(\cA^!))
  \simeq \cA$.
\item \emph{Interaction.}
  The two involutions are compatible: they are the
  color projections of the single $\SCchtop$-algebra
  quasi-isomorphism
  $\Omega^{\SCchtop}(\barB^{\SCchtop}(\cA^!))
  \simeq \cA$.
\end{enumerate}
\end{theorem}

\begin{proof}
The homotopy-Koszulity of $\SCchtop$
(Theorem~\ref{thm:homotopy-Koszul}) provides a
bar-cobar Quillen equivalence for $\SCchtop$-algebras
(Loday--Vallette~\cite{LV12}, Theorem~10.1.13; generalized
to colored operads by Vallette~\cite{Val07}).
By the self-duality $(\SCchtop)^! \simeq \SCchtop$
(Theorem~\ref{thm:SC-self-duality}), the bar-cobar
composite applied to $\cA^!$ yields an $\SCchtop$-algebra
quasi-isomorphic to the $\SCchtop$-algebra that
produced~$\cA^!$ in the first place, namely~$\cA$.
The color projections~(i) and~(ii) follow from the
naturality of the bar-cobar adjunction with respect to
the color projection functors.  Part~(iii) is the
statement that the bar-cobar composite is a natural
transformation of functors on $\SCchtop$-algebras,
refining the separate color involutions into a single
two-colored identity.
\end{proof}

\begin{remark}[What the involution says about Yangians]
\label{rem:involution-yangians}
\index{Yangian!involution}
For $\cA = \widehat{\fg}_k$, the Koszul dual is
$\cA^! = Y_\hh(\fg)$ (the Yangian).  The involution gives
$(Y_\hh(\fg))^! \simeq \widehat{\fg}_k$ as
$\SCchtop$-algebras: the dg-shifted Yangian of the Yangian
is the original affine Kac--Moody algebra.  The two color
projections are:
\[
  \underbrace{
  \Omegach(\barBch(Y_\hh(\fg)))
  \simeq \widehat{\fg}_k
  }_{\text{chiral bar-cobar inversion}}
  \qquad
  \underbrace{
  \Omega^{\mathrm{ord}}(\barB^{\mathrm{ord}}(
  Y_\hh(\fg))) \simeq \widehat{\fg}_k
  }_{\text{$E_1$ bar-cobar inversion}}.
\]
The first recovers the chiral algebra from the Sklyanin
$\lambda$-bracket; the second recovers it from the
Yangian product.  The $\SCchtop$-involution refines both
into a single identity, ensuring that the chiral and
$E_1$ recoveries are compatible via the mixed-color
operations.
\end{remark}

\subsection{Connection to Latyntsev's factorisation quantum groups}

The spectral $R$-matrix and the meromorphic tensor product are an instance of Latyntsev's \emph{factorisation quantum groups} \cite{Lat23}.

\begin{definition}[Factorisation Quantum Group (Informal)]
\label{def:factorisation_QG}
A \emph{factorisation quantum group} is a categorical structure encoding:
\begin{itemize}
\item A monoidal category $\mathcal{C}$ of representations;
\item A meromorphic tensor product $\otimes_z : \mathcal{C} \times \mathcal{C} \to \mathcal{C}[[z^{-1}, z]]$, parametrized by a spectral parameter $z \in \C$;
\item An $R$-matrix $R(z) : V \otimes_z W \to W \otimes_{-z} V$ satisfying the quantum Yang--Baxter equation.
\end{itemize}
\end{definition}

\begin{theorem}[Line Operators Form a Factorisation Quantum Group;
\ClaimStatusProvedHere]
\label{thm:lines_factorisationQG}
Let\/ $\cA$ be a logarithmic\/ $\SCchtop$-algebra.
The category $\mathcal{C}_{\text{line}}$ of line operators in a 3d HT theory, equipped with the meromorphic OPE and spectral $R$-matrix $r(z)$, forms a factorisation quantum group in the sense of Latyntsev \cite{Lat23}.
\end{theorem}

\begin{proof}
We verify the three axioms of a factorisation quantum group (Definition~\ref{def:factorisation_QG}).

\medskip
\noindent\textbf{(i) Monoidal structure.}
The category $\mathcal{C}_{\mathrm{line}}$ has objects given by boundary line operators and morphisms given by boundary-localized local operators (intertwiners).  The monoidal product is concatenation in the $E_1$-direction along the boundary $t = 0$: given lines $L_1, L_2$ placed at distinct points in $\C$, the product $L_1 \otimes L_2$ is defined by the $W(\SCchtop)$-module structure (Section~\ref{subsec:boundary-module}).  Associativity and unit constraints follow from the $E_1$-coherence encoded in the $W$-resolution.  By Theorem~\ref{thm:lines_as_modules}, $\mathcal{C}_{\mathrm{line}}$ is equivalent (as a dg category) to the category of $\A^!$-modules.

\medskip
\noindent\textbf{(ii) Meromorphic tensor product.}
For each $z \in \C^\times$, the OPE of line operators at separation $z$ defines a meromorphic family of tensor functors
\[
\otimes_z \colon \mathcal{C}_{\mathrm{line}} \times \mathcal{C}_{\mathrm{line}} \longrightarrow \mathcal{C}_{\mathrm{line}}[[z^{-1}, z]].
\]
Concretely, $L_1 \otimes_z L_2$ is the line operator obtained by placing $L_1$ at position $z$ and $L_2$ at position $0$ and taking the bulk-to-boundary OPE.  The OPE coefficients $C_k(z)$ are meromorphic in $z$ with poles only at $z = 0$ because the defining kernels are logarithmic along the collision divisor, as supplied by Theorem~\ref{thm:FM-calculus}; in physical realizations this matches the propagator statement furnished by Theorem~\ref{thm:physics-bridge}.  The family $\otimes_z$ varies holomorphically in $z$ on $\C^\times$ for the same reason: away from the collision divisor, the FM-compactified configuration integrals are smooth in the spectral parameter.

\medskip
\noindent\textbf{(iii) Braiding by $R(z)$.}
The spectral braiding $R(z) \colon L_1 \otimes_z L_2 \xrightarrow{\sim} L_2 \otimes_{-z} L_1$ (Definition~\ref{def:spectral-braiding}) is an isomorphism between $\otimes_z$ and the opposite product $\otimes_{-z}^{\mathrm{op}}$.  By Theorem~\ref{thm:YBE}, $R(z)$ satisfies the quantum Yang--Baxter equation---the braid relation for a braided monoidal category with spectral parameter.

\medskip
\noindent\textbf{(iv) Factorisation property.}
The distinguishing feature of a \emph{factorisation} quantum group (as opposed to an ordinary quantum group) is that the braiding depends holomorphically on $z$, upgrading the braided monoidal category to a \emph{factorisation category} on $\C$.  In Latyntsev's framework \cite{Lat23}, this means: the assignment $z \mapsto (\mathcal{C}_{\mathrm{line}}, \otimes_z, R(z))$ defines a constructible cosheaf on the Ran space $\Ran(\C)$.  For our category, this factorisation structure is inherited from the prefactorization algebra $\Obs^{\partial}$ on $\C \times \R_{\geq 0}$: restricting $\Obs^{\partial}$ to disks at different points of $\C$ and taking the $E_1$-coinvariants in the $\R$-direction yields fibers of the cosheaf, and the factorisation isomorphisms come from the prefactorization algebra structure maps for disjoint disks.

The holomorphic variation of the braiding in $z$ is therefore a direct consequence of the holomorphic-topological factorisation algebra structure of the 3d HT theory, restricted to the boundary.
\end{proof}

\subsection{Example: abelian CS and Heisenberg Yangian}

\begin{example}[Heisenberg Yangian from $U(1)$ CS]
\label{ex:Heisenberg_Yangian}
For abelian $U(1)$ Chern--Simons at level~$k$ (Example~\ref{ex:abelian_CS_complete}), the bulk algebra is the affine Kac--Moody $\widehat{\mathfrak{u}(1)}_k$; the Koszul dual is the Heisenberg Yangian $Y(\mathfrak{u}(1))$.  The spectral $R$-matrix is
\begin{equation}
r(z) = k \, \frac{J \otimes J}{z} + \text{(higher terms)},
\end{equation}
satisfying the classical Yang--Baxter equation.  For non-abelian $G$, the Koszul dual becomes the Yangian $Y(\mathfrak{g})$ or quantum affine algebra $U_q(\widehat{\mathfrak{g}})$, depending on the CS level and boundary conditions.
\end{example}

%%%%%%%%%%%%%%%%%%%%%%%%%%%%%%%%%%%%%%%%%%%%%%%%%%%%%%%%%%%%%%%%%%%%%%%%%%%%%
% THE E_1-CHIRAL STRUCTURE OF THE SHIFTED YANGIAN
%%%%%%%%%%%%%%%%%%%%%%%%%%%%%%%%%%%%%%%%%%%%%%%%%%%%%%%%%%%%%%%%%%%%%%%%%%%%%

\subsection{The $E_1$-chiral structure of the shifted Yangian}
\label{subsec:E1-chiral-yangian}
\index{dg-shifted Yangian!$E_1$-chiral structure|textbf}
\index{Yangian!spectral OPE}
\index{Yangian!PVA descent}
\index{Sklyanin bracket|textbf}

The Koszul dual $\cA^!$ is a dg-shifted Yangian
(Theorem~\ref{thm:Koszul_dual_Yangian}) and line operators
are $\cA^!$-modules (Theorem~\ref{thm:lines_as_modules}).
But $\cA^!$ carries a richer structure: it is an $E_1$-chiral
algebra on $\C_u \times \R_t$, with the $E_1$ direction from
the deconcatenation coproduct on $\R_t$ and the chiral direction
from the spectral OPE on $\C_u$.  We develop this structure
from the RTT relation, working from explicit formulas; the
definitive operadic framework is the master two-color Koszul
duality theorem (Theorem~\ref{thm:two-color-master}).

\subsubsection*{The FRT spectral OPE}

The starting point is the FRT commutation relation, derived
from the RTT relation $R_{12}(u{-}v)\,T_1(u)\,T_2(v) =
T_2(v)\,T_1(u)\,R_{12}(u{-}v)$
(equation~\eqref{eq:rtt-from-bar}).  For $\fg$ simple with
rational $R$-matrix $R(u) = 1 + \hbar\,P/u$ ($P$ the
permutation on $V \otimes V$, $V$ the defining representation),
the component form is
\begin{equation}\label{eq:FRT-components}
\boxed{
[t_{ij}(u),\, t_{kl}(v)]
\;=\;
\frac{\hbar}{u - v}\bigl(
  t_{kj}(u)\,t_{il}(v)
  \;-\;
  t_{kj}(v)\,t_{il}(u)
\bigr),
}
\end{equation}
where $T(u) = \sum_{I,J} E_{IJ} \otimes t_{IJ}(u) \in
\operatorname{End}(V) \otimes Y_\hbar(\fg)[\![u^{-1}]\!]$.
This is the \emph{spectral OPE}: it has the structure of a
chiral bracket with singularity at $u = v$.  The right-hand
side is a \emph{difference} of products evaluated at $u$ and
$v$, and vanishes at $u = v$---so the pole $1/(u-v)$ is
cancelled, making the bracket well-defined as a formal
distribution.

\subsubsection*{Calibration: the Heisenberg case}

\begin{computation}[Heisenberg Koszul dual as $E_1$-chiral
algebra; \ClaimStatusProvedHere]
\label{comp:heisenberg-E1-chiral}
For $\cA = \widehat{\mathfrak{u}(1)}_k$, the Koszul dual is
$\cA^! = \widehat{\mathfrak{u}(1)}_{-k}$
(Vol~I, Theorem~B).  There is a single FRT generator $t(u)$,
and the $R$-matrix is $R(u) = 1 + k/u$ (scalar, since
$\dim V = 1$).  The FRT relation~\eqref{eq:FRT-components}
gives
\[
[t(u),\, t(v)]
\;=\;
\frac{k}{u - v}\bigl(
  t(u)\,t(v) - t(v)\,t(u)
\bigr)
\;=\;
\frac{k}{u - v}\,[t(u), t(v)].
\]
This forces $[t(u), t(v)] = 0$ for $k \ne 0$
\textup(the zero mode is abelian\textup).
Passing to the current $J(u) = \partial_u \log t(u)$: the
relation becomes
\begin{equation}\label{eq:heisenberg-spectral-OPE}
[J_m,\, J_n] \;=\; -k\,m\,\delta_{m+n,0},
\end{equation}
and the spectral OPE is
$J(u)\,J(v) \sim -k/(u-v)^2$.
This is the Heisenberg algebra at level $-k$: the Koszul dual
flips the sign of the level.

The $E_1$-chiral structure is manifest: the spectral OPE has a
double pole at $u = v$ (a ``chiral'' singularity in the spectral
variable), and the algebra on $\R_t$ is the ordinary
Heisenberg---associative with no higher homotopies.  The
$A_\infty$ operations $m_k = 0$ for $k \ge 3$: the Heisenberg
Yangian is \emph{formal}.  This is the abelian calibration.
\end{computation}

\subsubsection*{The $\mathfrak{sl}_2$ computation}

For $\fg = \mathfrak{sl}_2$, the transfer matrix is
$T(u) = \bigl(\begin{smallmatrix}
t_{11}(u) & t_{12}(u) \\ t_{21}(u) & t_{22}(u)
\end{smallmatrix}\bigr)$
with $R(u) = 1 + \hbar P/u$,
$P = \sum_{i,j} E_{ij} \otimes E_{ji}$.
Equation~\eqref{eq:FRT-components} gives $16$ relations, of
which the independent ones are \textup(writing
$w = u - v$\textup):
\begin{align}
[t_{11}(u),\, t_{21}(v)]
&\;=\;
\frac{\hbar}{w}\bigl(t_{21}(u)\,t_{11}(v)
  - t_{21}(v)\,t_{11}(u)\bigr),
\label{eq:sl2-RTT-1}
\\[3pt]
[t_{21}(u),\, t_{12}(v)]
&\;=\;
\frac{\hbar}{w}\bigl(t_{11}(u)\,t_{22}(v)
  - t_{11}(v)\,t_{22}(u)\bigr),
\label{eq:sl2-RTT-2}
\\[3pt]
[t_{21}(u),\, t_{21}(v)]
&\;=\;
\frac{\hbar}{w}\bigl(t_{21}(u)\,t_{21}(v)
  - t_{21}(v)\,t_{21}(u)\bigr)
\;=\; 0.
\label{eq:sl2-RTT-3}
\end{align}
Relation~\eqref{eq:sl2-RTT-2} is the crucial one: it produces
the $[E, F] \sim H$ commutation relation of the Yangian.

In the Gauss decomposition
$T(u) = \bigl(\begin{smallmatrix} 1 & 0 \\ E(u) & 1
\end{smallmatrix}\bigr)
\bigl(\begin{smallmatrix} k_1(u) & 0 \\ 0 & k_2(u)
\end{smallmatrix}\bigr)
\bigl(\begin{smallmatrix} 1 & F(u) \\ 0 & 1
\end{smallmatrix}\bigr)$,
the Drinfeld generators $(E(u), F(u), H(u) = k_1(u)/k_2(u))$
satisfy
\begin{align}
(u - v - \hbar)\,E(u)\,E(v)
&\;=\;
(u - v + \hbar)\,E(v)\,E(u),
\label{eq:drinfeld-EE}
\\[3pt]
[E(u),\, F(v)]
&\;=\;
\frac{\hbar}{u - v}
\bigl(k_1(u)\,k_2(v) - k_1(v)\,k_2(u)\bigr),
\label{eq:drinfeld-EF}
\end{align}
derived from~\eqref{eq:sl2-RTT-1}--\eqref{eq:sl2-RTT-2} by
substitution.  These are the \emph{exact} spectral OPE
relations---not truncated at leading order.

\begin{remark}[What the spectral OPE is and is not]
\label{rem:spectral-OPE-nature}
Equations~\eqref{eq:FRT-components}
and~\eqref{eq:drinfeld-EE}--\eqref{eq:drinfeld-EF} are
\emph{not} $\lambda$-brackets in the PVA sense.  They are
commutation relations between generating functions, with
singularity at $u = v$.  The term ``spectral OPE'' is used
by analogy: the pole structure at $u = v$ parallels the pole
structure of a chiral algebra OPE at $z = w$, and the
residues encode the ``$E_1$-chiral product'' of line operators
at coincident spectral position.  But the generating functions
$E(u)$, $F(u)$, $H(u)$ are \emph{not} conformal fields: they
have no conformal weight, no state-field correspondence, and
no normally-ordered product.  The $E_1$-chiral structure on the
Yangian is a \emph{filtered} chiral algebra in the spectral
variable, not a vertex algebra.
\end{remark}

\subsubsection*{The classical limit: Sklyanin--Drinfeld Poisson
bracket}

\begin{proposition}[Sklyanin bracket from classical RTT;
\ClaimStatusProvedHere]
\label{prop:sklyanin-from-RTT}
\index{Sklyanin bracket|textbf}
The classical limit $\hbar \to 0$ of the FRT
relation~\eqref{eq:FRT-components}, with the replacement
$[\cdot, \cdot]/\hbar \to \{\cdot, \cdot\}$, gives
\begin{equation}\label{eq:sklyanin-bracket}
\{t_{ij}(u),\, t_{kl}(v)\}
\;=\;
\frac{1}{u - v}\bigl(
  t_{kj}(u)\,t_{il}(v) - t_{kj}(v)\,t_{il}(u)
\bigr).
\end{equation}
This is the \emph{Sklyanin--Drinfeld Poisson bracket} on the
classical Yangian $Y_{\mathrm{cl}}(\fg) \simeq
\mathcal{O}(G[\![u^{-1}]\!])$, the coordinate ring of the
formal arc group.
\end{proposition}

\begin{proof}
Divide~\eqref{eq:FRT-components} by $\hbar$ and take
$\hbar \to 0$.  The left side
$[t_{ij}(u), t_{kl}(v)]/\hbar$ becomes the Poisson bracket
$\{t_{ij}(u), t_{kl}(v)\}$.  The right side
$(t_{kj}(u)\,t_{il}(v) - t_{kj}(v)\,t_{il}(u))/(u-v)$
is already independent of $\hbar$ (the generators $t_{ij}(u)$
have classical limit $t_{ij}^{\mathrm{cl}}(u)$).
\end{proof}

\begin{proposition}[Sklyanin bracket satisfies Jacobi;
\ClaimStatusProvedHere]
\label{prop:sklyanin-jacobi}
The bracket~\eqref{eq:sklyanin-bracket} satisfies the Poisson
Jacobi identity.
\end{proposition}

\begin{proof}
The Jacobi identity for~\eqref{eq:sklyanin-bracket} is
equivalent to the classical Yang--Baxter equation (CYBE) for
$r(u) = P/u$:
$[r_{12}(u{-}v),\, r_{13}(u{-}w)]
+ [r_{12}(u{-}v),\, r_{23}(v{-}w)]
+ [r_{13}(u{-}w),\, r_{23}(v{-}w)] = 0$.
This is verified by direct computation: each commutator
produces terms of the form $P_{ij}P_{kl}/(u_a - u_b)(u_c - u_d)$,
and the six terms cancel pairwise by the permutation identity
$P_{12}P_{13} + P_{13}P_{23} + P_{23}P_{12}
= P_{12}P_{23} + P_{23}P_{13} + P_{13}P_{12}$.
Alternatively: this is the classical limit of the quantum
Yang--Baxter equation (Theorem~\ref{thm:YBE}), which is
proved in~\S\ref{subsec:YBE}.
\end{proof}

\begin{computation}[Heisenberg verification;
\ClaimStatusProvedHere]
\label{comp:heisenberg-sklyanin}
For $\fg = \mathfrak{u}(1)$: $t(u) = J(u)$,
$P = 1$, and
$\{J(u), J(v)\} = (J(u) J(v) - J(v) J(u))/(u-v) = 0$
(the classical Heisenberg is abelian).  The Poisson bracket on
modes $J_r$ comes from the second-order pole of the
\emph{quantum} relation~\eqref{eq:heisenberg-spectral-OPE}:
$\{J_m, J_n\}_{\mathrm{qu}} = -k\,m\,\delta_{m+n,0}$.
This is the standard Heisenberg PVA at level $-k$, confirming
that $\cA^! = \widehat{\mathfrak{u}(1)}_{-k}$.
\end{computation}

\subsubsection*{The $A_\infty$ structure: what is proved and
what is open}

The $E_1$-chiral $A_\infty$ structure on $Y_\hbar(\fg)$ is
the central open question.  We state precisely what is known.

The Yangian $Y_\hbar(\fg)$ is a \emph{filtered} algebra:
the PBW filtration by total mode number $|t_{ij}^{(r)}| = r$
gives $\gr Y_\hbar(\fg) \simeq U(\fg[t])$.

\begin{proposition}[Associated graded is formal;
\ClaimStatusProvedHere]
\label{prop:gr-yangian-formal}
The associated graded $U(\fg[t])$ is Koszul as a graded
algebra.  Its bar complex is generated in degrees $1$ and $2$:
the generators are $\fg[t]$ in degree~$1$, and the relations
$xy - yx - [x,y] = 0$ in degree~$2$.  The $A_\infty$ structure
on $U(\fg[t])$ is formal: $m_k = 0$ for $k \ge 3$.
\end{proposition}

\begin{proof}
$U(\fg[t]) = U(\fg) \otimes \C[t]$ as a filtered algebra.
Both $U(\fg)$ (Priddy) and $\C[t]$ (trivially) are Koszul, and
the tensor product of Koszul algebras is Koszul
(Polishchuk--Positselski, Theorem~2.3.1).
\end{proof}

\begin{question}[$E_1$-chiral $A_\infty$ of $Y_\hbar(\fg)$;
\ClaimStatusOpen]
\label{q:yangian-Ainfty}
\index{Yangian!$A_\infty$ structure}
Is the full $Y_\hbar(\fg)$ (not just its associated graded)
formal as an $E_1$-chiral algebra?  The RTT
relation~\eqref{eq:FRT-components} is quadratic in the
generating functions $t_{ij}(u)$, but the mode-expanded
relations~\eqref{eq:yangian-relation-modes} mix all mode
numbers.  The algebra is a non-trivial filtered deformation
of a Koszul algebra.

Two scenarios:
\begin{enumerate}[label=\textup{(\alph*)}]
\item \emph{Filtered formality:}
  There exists a filtered $A_\infty$ quasi-isomorphism
  $Y_\hbar(\fg) \xrightarrow{\sim} (U(\fg[t]), m_2)$
  that is the identity on the associated graded.  This would
  mean $Y_\hbar(\fg)$ is ``formal up to filtration'': the
  higher $m_k$ vanish on $\gr$ but may be non-zero on the
  filtered level.
\item \emph{Genuine non-formality:}
  No such quasi-isomorphism exists, and the $A_\infty$
  operations $m_k$ for $k \ge 3$ are genuinely non-trivial
  (not killed by any change of generators).
\end{enumerate}
The answer depends on the vanishing of specific deformation
classes in the Harrison cohomology of $U(\fg[t])$ with
coefficients in the adjoint bimodule, twisted by the
$\hbar$-deformation.  This is not computed in either volume.
\end{question}

\begin{remark}[Virasoro: non-formality on both sides]
\label{rem:yangian-virasoro-nonformality}
For the Virasoro algebra $\mathrm{Vir}_c$ at generic $c$,
the Koszul dual $\mathrm{Vir}_{26-c}$ is \emph{not} formal:
$m_k \ne 0$ for all $k \ge 3$
(Section~\ref{subsec:gravity-shadow-tower}).  The Virasoro
has no Yangian simplification: both sides of the Koszul
duality are genuinely non-formal, and the infinite-depth
shadow tower persists on the line-operator side.
\end{remark}

\subsubsection*{The bar-cobar involution}

The Koszul involution $(-)^{!!} \simeq \id$ at the
$E_1$-chiral level gives
$\Omega^{E_1\text{-ch}}(B^{E_1\text{-ch}}(\cA)) \simeq \cA$.
For $\cA = \widehat{\fg}_k$: the ordered bar produces
$Y_\hbar(\fg)$, and the cobar recovers $\widehat{\fg}_k$.

The structures exchange:
\begin{center}
\small
\renewcommand{\arraystretch}{1.15}
\begin{tabular}{lll}
\textbf{Boundary $\cA$} &
  \textbf{Koszul dual $\cA^!$} &
  \textbf{Mechanism} \\
\hline
Current OPE $J^a(z)J^b(w)$ &
  FRT relation~\eqref{eq:FRT-components} &
  bar $\leftrightarrow$ cobar \\[2pt]
Poles in $z - w$ &
  Poles in $u - v$ &
  holomorphic $\leftrightarrow$ spectral \\[2pt]
$\lambda$-bracket &
  Sklyanin bracket~\eqref{eq:sklyanin-bracket} &
  PVA $\leftrightarrow$ PVA \\[2pt]
$r(z) = \Omega/z$ &
  $r_{\mathrm{cl}}(u) = P/u$ &
  same Casimir \\[2pt]
$\kappa = \frac{\dim\fg\,(k+h^\vee)}{2h^\vee}$ &
  $\kappa^! = -\frac{\dim\fg\,(k+h^\vee)}{2h^\vee}$ &
  complementarity \\[2pt]
$\SCchtop$-algebra &
  $\SCchtop$-algebra &
  two-color involution (Thm~\ref{thm:two-color-master})
\end{tabular}
\end{center}

\begin{computation}[Heisenberg bar-cobar;
\ClaimStatusProvedHere]
\label{comp:heisenberg-bar-cobar}
For $\cA = \widehat{\mathfrak{u}(1)}_k$: the boundary has
$J(z)J(w) \sim k/(z-w)^2$, and $\cA^!$ has
$J(u)J(v) \sim -k/(u-v)^2$.  The bar complex of $\cA^!$
should recover $\cA$.  Indeed: the ordered bar differential on
$[s^{-1}J \,|\, s^{-1}J]$ extracts the
residue of $-k/(u-v)^2$, which is $k\,\partial$---the level-$k$
Heisenberg OPE.  The level flips sign twice under the double
bar and returns to $k$: $(\cA^!)^! = \cA$.
\end{computation}

\subsubsection*{Genus tower from complementarity}

Koszul complementarity (Vol~I, Theorem~C) gives
$\kappa(\cA) + \kappa(\cA^!) = K(\cA)$ in general, where
$K(\cA)$ is a level-independent complementarity constant.
For the affine lineage ($K = 0$):
\begin{equation}\label{eq:yangian-kappa}
\kappa(\cA^!)
\;=\;
-\kappa(\cA).
\end{equation}
For $\cA = \widehat{\fg}_k$: $\kappa = \dim\fg\,(k+h^\vee)/(2h^\vee)$, so
$\kappa(Y_\hbar(\fg)) = -\dim\fg\,(k+h^\vee)/(2h^\vee)$.  The complementarity
potential $F_1(\cA) + F_1(\cA^!) = 0$ vanishes: the bulk
theory carries no net genus-$1$ anomaly.

The all-genus generating function:
\[
\sum_{g \ge 1} F_g(Y_\hbar(\fg))\,\hbar^{2g-2}
\;=\;
-\frac{k\dim\fg}{2}\bigl(\hat{A}(i\hbar) - 1\bigr).
\]

\begin{computation}[Heisenberg genus tower;
\ClaimStatusProvedHere]
\label{comp:heisenberg-genus}
For $\cA = \widehat{\mathfrak{u}(1)}_k$:
$\kappa = k$, $\kappa^! = -k$.  The genus-$1$ free energies
are $F_1(\cA) = k/24$ and $F_1(\cA^!) = -k/24$.  Their sum
vanishes: the bulk $\mathfrak{u}(1)$ theory has no
genus-$1$ anomaly.  This matches the direct computation:
the bulk partition function on $T^2$ is
$Z_{\mathrm{bulk}} = |\eta(\tau)|^{-2}$, which is modular
invariant for all $k$.
\end{computation}

\begin{remark}[What remains open]
\label{rem:yangian-E1-chiral-open}
Four problems are open:
\begin{enumerate}[label=\textup{(\roman*)}]
\item The $A_\infty$ formality question
  (Question~\ref{q:yangian-Ainfty}).
\item Construction of a ``Yangian vertex algebra''---a vertex
  algebra-like object on $Y_\hbar(\fg)$ encoding the spectral
  OPE, state-field correspondence, and translation.  The
  obstruction: the Yangian is filtered by mode number, not
  graded by conformal weight.  Theorem~\ref{thm:dual-sc-algebra}
  resolves the \emph{structural} question---$\cA^!$ carries a
  canonical $\SCchtop$-algebra structure encoding both the
  Sklyanin $\lambda$-bracket (closed color) and the Yangian
  product (open color)---while the explicit ``vertex algebra''
  presentation remains a presentation problem.
\item Quantization of the Sklyanin PVA
  \eqref{eq:sklyanin-bracket} via the modular programme
  (Chapter~\ref{ch:modular-pva-quantization}).
\item The $E_1$-chiral $A_\infty$ structure on
  non-formal Koszul duals ($\mathrm{Vir}_{26-c}$,
  $\mathcal{W}$-algebras at dual parameters).
\end{enumerate}
\end{remark}


%%%%%%%%%%%%%%%%%%%%%%%%%%%%%%%%%%%%%%%%%%%%%%%%%%%%%%%%%%%%%%%%%%%%%%%%%%%%%
% CURVATURE-BRAIDING DECOUPLING
%%%%%%%%%%%%%%%%%%%%%%%%%%%%%%%%%%%%%%%%%%%%%%%%%%%%%%%%%%%%%%%%%%%%%%%%%%%%%

\subsection{Curvature-braiding decoupling: the Steinberg perspective}
\label{subsec:curvature-braiding-decoupling}
\index{curvature-braiding decoupling|textbf}
\index{Steinberg variety!curvature-braiding decoupling}
\index{genus tower!curvature-braiding dichotomy}

\providecommand{\Mvac}{\mathcal{M}_{\mathrm{vac}}}
\providecommand{\Steinb}{\mathfrak{S}}

The genus-$g$ bar complex carries a curvature $\kappa(\cA)\cdot\omega_g$
from the Hodge bundle (Vol~I, Theorem~D) and a spectral $R$-matrix
$R(z)$ from the braiding of line operators
(\S\ref{subsec:line-ops}).  A central question---item~(iv) of
Remark~\ref{rem:yangian-E1-chiral-open}---is whether these two structures
interact.  There are two potential channels of interaction, controlled
by different algebraic data:
\begin{enumerate}[label=(\alph*)]
\item \label{item:bar-channel}
  The \emph{bar-complex channel}: the curvature
  $\kappa\cdot\omega_g$ modifies the differential $d_g = d_0 +
  \kappa\cdot\omega_g$, which could affect the braiding via the
  compatibility of $d_g$ with the $E_1$-coproduct $\Delta$.
\item \label{item:propagator-channel}
  The \emph{propagator channel}: the genus-$g$ propagator on
  $\Sigma_g$ differs from the genus-$0$ propagator on $\C$; since
  the spectral $R$-matrix is computed from an exchange integral
  involving the holomorphic propagator
  (Proposition~\ref{prop:field-theory-r}), the $R$-matrix can
  acquire genus dependence through the propagator even when the
  bar-complex channel is trivial.
\end{enumerate}
For Koszul algebras, channel~\ref{item:bar-channel} is trivial
(Proposition~\ref{prop:curvature-braiding-decoupling} below).
Channel~\ref{item:propagator-channel} is controlled by the Coisson
bracket $c_0 = \{a {}_0 b\}$: it is trivial if and only if $c_0 = 0$
(Theorem~\ref{thm:elliptic-spectral-dichotomy}).  The full
decoupling criterion is therefore $c_0 = 0$, not Koszulness alone.

\subsubsection*{Bar-complex decoupling for Koszul algebras}

\begin{proposition}[Bar-complex curvature-braiding decoupling;
\ClaimStatusProvedHere]
\label{prop:curvature-braiding-decoupling}
\index{curvature-braiding decoupling!Koszul case}
Let $\cA$ be a chirally Koszul algebra.  Then:
\begin{enumerate}[label=\textup{(\roman*)}]
\item The reduced bar cohomology $H^2(\bar{B}(\cA))$ is concentrated
  in bar degree~$1$.  In particular, the curvature
  $\kappa(\cA) \in H^2(\bar{B}(\cA))$ is \emph{primitive}:
  it lies in the space of bar-degree-$1$ cocycles.
\item The genus-$g$ deformation $\kappa\cdot\omega_g$ is compatible
  with the $E_1$-coproduct $\Delta$: it acts as a coderivation
  on the bar coalgebra.
\item The $E_1$-coproduct structure on line operators is
  genus-independent: the coproduct $\Delta_g$ at genus $g$ is
  identical to $\Delta_0$.
\end{enumerate}
In particular, the bar-complex
channel~\textup{\ref{item:bar-channel}} contributes no genus
correction to the spectral $R$-matrix.

\medskip\noindent
\textbf{Warning.}  This does \emph{not} imply $R_g(z) = R_0(z)$
in general, because the spectral $R$-matrix also depends on the
holomorphic propagator via the exchange integral
\textup{(}Proposition~\textup{\ref{prop:field-theory-r})}.  At
genus~$g \ge 1$, the propagator acquires quasi-periodic monodromy
from $\FM_k(\Sigma_g)$, which can independently modify $R(z)$
through the propagator
channel~\textup{\ref{item:propagator-channel}}.  The full
decoupling criterion is the vanishing of the Coisson bracket
$c_0 = \{a {}_0 b\} = 0$
\textup{(}Theorem~\textup{\ref{thm:elliptic-spectral-dichotomy})},
which controls the propagator channel.  For abelian algebras
\textup{(}$c_0 = 0$\textup{)}, both channels are trivial and
$R_g(z) = R_0(z)$.  For non-abelian Koszul algebras
\textup{(}$c_0 \ne 0$, e.g., $\hat{\fg}_k$\textup{)}, the
bar-complex channel is trivial but the propagator channel is
not: $r^{E_\tau}(z) \ne r^{\C}(z)$
\textup{(}Theorem~\textup{\ref{thm:elliptic-spectral-dichotomy}},
item~\textup{(\ref{item:dichotomy-entanglement}))}.
\end{proposition}

\begin{proof}
By the Koszul hypothesis, the bar complex $\barB(\cA)$ is formal:
its cohomology is generated in bar degree~$1$ and the differential
$d_{\barB}$ raises bar degree by exactly one.  The curvature
$\kappa(\cA)$ lives in $H^2(\bar{B}(\cA))$, which is therefore
concentrated in bar degree~$1$.  (Note: $H^2$ in bar degree~$1$
may be multi-dimensional if $\cA$ has multiple independent
relations, but every element is primitive.)

The key point is that a bar-degree-$1$ curvature element is
a \emph{coderivation} of the bar coalgebra: it satisfies
$\Delta \circ \kappa = (\kappa \otimes \mathrm{id}
+ \mathrm{id} \otimes \kappa) \circ \Delta$.  This is
because bar-degree-$1$ elements correspond to generators of
the bar construction, on which the coproduct $\Delta$ is
primitive (deconcatenation of length-$1$ words).

The genus-$g$ deformation of the $\SCchtop$-algebra structure
adds $\kappa\cdot\omega_g$ to the bar differential.
Since $\kappa\cdot\omega_g$ is a coderivation, the deformed
differential $d_g = d_0 + \kappa\cdot\omega_g$ remains
compatible with~$\Delta$.  This establishes that the
$E_1$-coproduct structure is undeformed by the genus-$g$
curvature, closing the bar-complex channel.
\end{proof}

\subsubsection*{Entanglement for algebras with non-formal Swiss-cheese structure}

\begin{proposition}[Curvature-braiding entanglement;
\ClaimStatusProvedHere]
\label{prop:curvature-braiding-entanglement}
\index{curvature-braiding decoupling!non-formal case}
Let $\cA$ be a chiral algebra with non-formal Swiss-cheese structure
(e.g., $\mathrm{Vir}_c$ at generic~$c$).  Then:
\begin{enumerate}[label=\textup{(\roman*)}]
\item $H^2(\bar{B}(\cA))$ is multi-dimensional: the curvature
  $\kappa(\cA)$ lives in a space of $\dim > 1$.
\item Different line operators see different components of~$\kappa$.
\item The genus-$g$ $R$-matrix acquires corrections:
  \begin{equation}\label{eq:R-genus-corrections}
    R_g(z)
    \;=\;
    R_0(z)
    \;+\;
    \sum_{k \ge 1} R_k(z)\cdot(\kappa\cdot\omega_g)^k,
  \end{equation}
  where $R_k(z) \in \End(L_1 \otimes L_2)[\![z^{-1}]\!]$ are
  explicit corrections determined by the higher bar cohomology.
\end{enumerate}
\end{proposition}

\begin{proof}
When $\cA$ has non-formal Swiss-cheese structure, $\barB(\cA)$ has cohomology in bar
degrees $> 1$.  In particular $H^2(\bar{B}(\cA))$ receives
contributions from bar degrees $1, 2, \ldots$, making the
curvature vector-valued: $\kappa \in H^2(\bar{B}(\cA))$ with
$\dim H^2 > 1$.

When $H^2$ has components in bar degree $\ge 2$, those
components are \emph{not} coderivations of the bar coalgebra:
they do not satisfy the primitivity condition
$\Delta \circ \kappa = (\kappa \otimes \mathrm{id}
+ \mathrm{id} \otimes \kappa) \circ \Delta$.
Consequently, the deformed differential $d_g$ is no longer
compatible with~$\Delta$, and different line operators $L_i$
couple to different components $\kappa^{(j)}$ via the action
of $H^2(\bar{B}(\cA))$ on
$C_{\mathrm{line}} \simeq \cA^!\text{-}\mathsf{mod}$.
The genus-$g$ deformation $d_g = d_0 + \kappa\cdot\omega_g$
therefore modifies the braiding transport, and a standard
perturbative expansion in $\kappa\cdot\omega_g$ yields
\eqref{eq:R-genus-corrections}.  The $k$-th correction
$R_k(z)$ is determined by the bar $A_\infty$ operations
$m_j$ with $j \le k+1$ acting on $H^*(\bar{B}(\cA))$.
\end{proof}

\subsubsection*{The dichotomy theorem}

\begin{theorem}[Pole-structure dichotomy;
\ClaimStatusProvedHere]
\label{thm:pole-structure-dichotomy}
\index{curvature-braiding decoupling!pole-structure dichotomy}
\index{OPE poles!curvature-braiding dichotomy}
Let $\cA$ be a chirally Koszul logarithmic $\SCchtop$-algebra.
Consider the following conditions:
\begin{enumerate}[label=\textup{(\roman*)}]
\item All OPE poles of\/ $\cA$ are simple.
\item $\cA$ is chirally Koszul.
\item $H^2(\bar{B}(\cA))$ is concentrated in bar degree~$1$
  \textup{(}bar-complex decoupling\textup{)}.
\item The Coisson bracket vanishes: $c_0 = \{a {}_0 b\} = 0$
  for all $a,b \in H^\bullet(\cA,Q)$.
\item Full curvature-braiding decoupling holds: $R_g(z) = R_0(z)$
  for all $g \ge 0$.
\end{enumerate}
Then
$\textup{(ii)} \Rightarrow \textup{(i)}$,\quad
$\textup{(ii)} \Rightarrow \textup{(iii)}$,\quad
$\textup{(iii)} \;\&\; \textup{(iv)}
  \Rightarrow \textup{(v)}$.

\noindent
Condition~\textup{(iii)} alone gives bar-complex decoupling
\textup{(}coproduct $\Delta$ is genus-independent\textup{)}, but
this is weaker than~\textup{(v)}: the spectral $R$-matrix also
depends on the holomorphic propagator, whose quasi-periodic
monodromy at genus~$g \ge 1$ is controlled by $c_0$
\textup{(}Theorem~\textup{\ref{thm:elliptic-spectral-dichotomy}}).

\noindent
Conversely, if\/ $\cA$ has higher-order OPE poles, then
bar-complex entanglement is generically present
\textup{(}Proposition~\textup{\ref{prop:curvature-braiding-entanglement})},
and if $c_0 \ne 0$, then propagator-channel entanglement is
also present.

\smallskip\noindent
\textbf{Note.}  The implication
\textup{(i)}~$\Rightarrow$~\textup{(ii)} requires an
additional hypothesis: simple poles guarantee that the bar
complex is generated in degree~$1$, but Koszulness further
requires that all relations are quadratic
\textup{(}generated in bar degree~$2$\textup{)}.
For the standard examples, observe the two-dimensional
classification: the Heisenberg algebra $\cH_k$ and
lattice VOAs $V_\Lambda$ have $c_0 = 0$
\textup{(}full decoupling\textup{)}; the $\beta\gamma$
system has $c_0 = \beta \otimes \gamma
- \gamma \otimes \beta \ne 0$ \textup{(}from
$\{\beta {}_0 \gamma\} = 1$, the simple-pole
OPE\textup{)}, and affine Kac--Moody $\hat{\fg}_k$
has $c_0 = \Omega \ne 0$ \textup{(}the
Casimir\textup{)}.  Both $\beta\gamma$ and
$\hat{\fg}_k$ exhibit propagator entanglement at
genus~$\ge 1$.
\end{theorem}

\begin{proof}
(ii)~$\Rightarrow$~(i): Koszulness means the bar complex is
generated in degree~$1$ with quadratic relations.  Generation
in degree~$1$ forces all OPE poles to be simple (Vol~I,
Theorem~A: bar generators correspond to OPE singular terms).

(ii)~$\Rightarrow$~(iii): For a Koszul algebra, the bar
complex is cogenerated in bar degree~$1$
(the desuspended generators $s^{-1}\bar{\cA}$), and all
bar cohomology sits on the expected diagonal
$H^{p,q}(\barB(\cA)) = 0$ for $q \neq 0$
(equation~\eqref{eq:koszul-concentration}).  The curvature
class $\kappa(\cA) \in H^2(\bar{B}(\cA))$ is therefore
represented by a cocycle supported in bar degree~$1$
(the cogenerating part of the cofree coalgebra), hence
it is primitive for the deconcatenation coproduct.
Condition~(iii) follows.

(iii)~\&~(iv)~$\Rightarrow$~(v):
Condition~(iii) ensures that the curvature acts as a
coderivation of the bar coalgebra, so the $E_1$-coproduct
$\Delta$ is undeformed at genus~$g$
(Proposition~\ref{prop:curvature-braiding-decoupling}).
Condition~(iv) ensures that the genus-$g$ propagator on
$\Sigma_g$ produces no quasi-periodic monodromy in the
exchange integral: when $c_0 = 0$, the elliptic
$r$-matrix is doubly periodic
(Theorem~\ref{thm:elliptic-spectral-dichotomy},
item~(\ref{item:dichotomy-decoupling})), hence
$R_g(z) = R_0(z)$.

For the converse direction: if $\cA$ has a pole of order
$\ge 2$, then $\barB(\cA)$ has generators in bar degree
$\ge 2$.  The resulting higher-bar-degree contributions to
$H^2(\bar{B}(\cA))$ produce a curvature that is not a
coderivation, and
Proposition~\ref{prop:curvature-braiding-entanglement}
gives bar-complex entanglement $R_g \ne R_0$ generically.
If $c_0 \ne 0$, the propagator channel independently
contributes $r^{E_\tau}(z) \ne r^{\C}(z)$
(Theorem~\ref{thm:elliptic-spectral-dichotomy},
item~(\ref{item:dichotomy-entanglement})).
\end{proof}

\begin{remark}[Steinberg interpretation]
\label{rem:steinberg-curvature-braiding}
\index{Steinberg variety!curvature interpretation}
The dichotomy admits a geometric reading via the Steinberg
analogy of \S\ref{subsec:two-color-koszul-duality}.
The genus-$g$ curvature deforms the ambient vacuum stack
$\Mvac$: it acquires a curved structure from the Hodge
bundle $\omega_g \in H^0(\bar{\mathcal{M}}_g, \omega)$.
The braiding is the monodromy of the flat connection on
the Steinberg variety~$\Steinb_b$.

The bar-complex channel decouples when the Hodge
deformation is \emph{central}---acting by scalars on the
Borel--Moore homology $H_*^{\mathrm{BM}}(\Steinb_b)$.
This centrality holds if and only if $\Steinb_b$ is
\emph{concentrated} (cohomology in a single degree),
which is the geometric shadow of Koszulness.  When
$\Steinb_b$ has cohomology in multiple degrees, the
Hodge deformation acts non-centrally and the monodromy
(braiding) is modified via the bar-complex channel.

The propagator channel provides a second, independent
source of genus dependence: the holomorphic propagator
on $\Sigma_g$ (a section of a nontrivial line bundle)
modifies the exchange integral defining $R(z)$.  This
channel is controlled by the Coisson bracket $c_0$,
not by the concentration of $\Steinb_b$.  Full decoupling
requires both: concentrated self-intersection (Koszulness)
\emph{and} vanishing Coisson bracket ($c_0 = 0$).
\end{remark}

\begin{remark}[Curvature-braiding as the perturbative/non-perturbative watershed]
\label{rem:curvature-braiding-perturbative}
\index{curvature-braiding decoupling!perturbative exactness}
\index{perturbative!curvature-braiding decoupling}
\index{non-perturbative!curvature-braiding entanglement}
Full curvature-braiding decoupling (both bar-complex and propagator
channels trivial) is perturbative exactness: the Lagrangian
self-intersection is clean, the genus tower acts by scalars, and the
$R$-matrix is determined once and for all at genus~$0$.  Partial
decoupling (bar-complex channel closed, propagator channel open)
occurs for non-abelian Koszul algebras: the coproduct is
genus-independent, but the $R$-matrix acquires quasi-periodic
corrections from the genus-$g$ propagator.  Full
entanglement (both channels active) occurs for algebras with
non-formal Swiss-cheese structure: the Lagrangian self-intersection
has excess, higher $\Ainf$ operations couple the curvature to the
braiding, and each genus level contributes genuinely new spectral
data.  The dichotomy is controlled by two independent invariants---the
Lagrangian intersection type (transverse vs.\ excess, i.e.\ formal
vs.\ non-formal Swiss-cheese) and the Coisson bracket
($c_0 = 0$ vs.\ $c_0 \ne 0$)---which refine the
perturbative/non-perturbative divide
(Remark~\ref{rem:koszul-perturbative-watershed}).
\end{remark}

\begin{corollary}[Genus corrections and formality failure;
\ClaimStatusProvedHere]
\label{cor:genus-corrections-formality}
\index{formality!failure at $d'=1$}
For a chirally Koszul algebra, the bar-complex channel is trivial:
$d^2 = \kappa\cdot\omega_g$ acts as a coderivation, and the
$E_1$-coproduct is genus-independent.  The remaining genus
dependence of the spectral $R$-matrix is controlled entirely by
the Coisson bracket $c_0$
\textup{(}Theorem~\textup{\ref{thm:elliptic-spectral-dichotomy})}:
if $c_0 = 0$ \textup{(}abelian case\textup{)}, no genus
corrections arise; if $c_0 \ne 0$
\textup{(}non-abelian case\textup{)}, the propagator channel
contributes quasi-periodic corrections at genus~$g \ge 1$.

For an algebra with non-formal Swiss-cheese structure, both channels produce corrections:
the bar-complex channel contributes $R_k(z)$ corrections from
the higher bar cohomology, and the propagator channel contributes
quasi-periodic modifications.  The failure of
formality at $d' = 1$ \textup{(}the non-vanishing of higher $A_\infty$
operations\textup{)} is precisely the manifestation of
bar-complex entanglement---the higher $m_k$ encode the
coupling between the multi-dimensional curvature and the spectral
braiding through the bar-complex channel.
\end{corollary}

\begin{proof}
For $\cA$ Koszul, Proposition~\ref{prop:curvature-braiding-decoupling}
gives bar-complex decoupling: the coproduct $\Delta$ is
genus-independent.  The remaining genus dependence enters through
the propagator channel, classified by
Theorem~\ref{thm:elliptic-spectral-dichotomy}.

When $\cA$ has non-formal Swiss-cheese structure, the higher $A_\infty$ operations
$m_k$ with $k \ge 3$ are non-trivial (formality fails).
By Theorem~\ref{thm:pole-structure-dichotomy}, the curvature
is vector-valued, and the corrections $R_k(z)$ in
\eqref{eq:R-genus-corrections} are determined by these $m_k$.
Thus formality failure and bar-complex entanglement
are two descriptions of the same phenomenon.
\end{proof}


%%%%%%%%%%%%%%%%%%%%%%%%%%%%%%%%%%%%%%%%%%%%%%%%%%%%%%%%%%%%%%%%%%%%%%%%%%%%%
% TWO-COLOR KOSZUL DUALITY
%%%%%%%%%%%%%%%%%%%%%%%%%%%%%%%%%%%%%%%%%%%%%%%%%%%%%%%%%%%%%%%%%%%%%%%%%%%%%

\subsection{Two-color Koszul duality}
\label{subsec:two-color-koszul-duality}
\index{two-color Koszul duality|textbf}
\index{Swiss-cheese operad!two-color Koszul duality}
\index{Koszul duality!two-color}

The bar complex $\barB(\cA)$ carries two compatible structures:
the bar differential $d_{\barB}$ (from OPE residues on
$\FM_k(\C)$) and the deconcatenation coproduct $\Delta$
(from ordered configurations on $\Conf_k^{<}(\R)$).
Together these make $\barB(\cA)$ an $\SCchtop$-coalgebra.
The homotopy-Koszulity of $\SCchtop$
(Theorem~\ref{thm:homotopy-Koszul}) and the associated graded
splitting (Proposition~\ref{prop:gr-chiral}) produce a
\emph{simultaneous} Koszul duality in both colors.  This
subsection extracts the master theorem.

\subsubsection*{The two-color Koszul datum}

\begin{definition}[Two-color Koszul datum]
\label{def:two-color-koszul-datum}
\index{two-color Koszul datum}
Let $\cA$ be a chirally Koszul logarithmic
$\SCchtop$-algebra and let $\cA^!$ be its Koszul dual.  The \emph{two-color Koszul
datum} is the sextuple
\[
\bigl(\cA,\; \cA^!,\;
\barBch(\cA),\; \barB^{\mathrm{ord}}(\cA),\;
\tau_{\mathrm{ch}},\; \tau_{\mathrm{ord}}\bigr),
\]
where:
\begin{enumerate}[label=\textup{(\roman*)}]
\item $\barBch(\cA)$ is the \emph{chiral bar complex}
  (Definition~\ref{def:chiral_bar}): the bar construction with
  differential from $\FM_k(\C)$-integrals.  This is the
  closed-color sector of the $\SCchtop$-bar complex.
\item $\barB^{\mathrm{ord}}(\cA)$ is the \emph{ordered bar
  complex}: the cofree coalgebra $T^c(s^{-1}\bar{\cA})$ with
  deconcatenation coproduct from $\Conf_k^{<}(\R)$.  This is
  the open-color sector.
\item $\tau_{\mathrm{ch}} \colon \barBch(\cA) \to \cA^!$ and
  $\tau_{\mathrm{ord}} \colon \barB^{\mathrm{ord}}(\cA) \to
  \cA^!$ are the closed-color and open-color universal twisting
  morphisms, respectively.
\end{enumerate}
The two bar complexes are \emph{not independent}: they are the
two projections of the single $\SCchtop$-bar complex
$\barB^{\SCchtop}(\cA)$.
\end{definition}

\subsubsection*{The master theorem}

\begin{theorem}[Two-color Koszul duality;
\ClaimStatusProvedHere]
\label{thm:two-color-master}
\index{two-color Koszul duality!master theorem|textbf}
Let $\cA$ be a chirally Koszul logarithmic
$\SCchtop$-algebra.  The homotopy-Koszulity of $\SCchtop$
simultaneously produces:
\begin{enumerate}[label=\textup{(\roman*)}]
\item \emph{Closed-color duality.}
  A Quillen equivalence between chiral $\cA$-modules and
  chiral $\cA^!$-comodules, recovering
  Theorems~A and~B of~Vol~I\@.  The bar functor
  $\barBch \colon \cA\text{-mod}_{\mathrm{ch}} \rightleftarrows
  \cA^!\text{-comod}_{\mathrm{ch}} \cocolon \Omegach$ is
  the closed-color bar-cobar adjunction.
\item \emph{Open-color duality.}
  A Quillen equivalence between $E_1$ $\cA$-modules and $E_1$
  $\cA^!$-comodules, identifying
  $\mathcal{C}_{\mathrm{line}} \simeq \cA^!\text{-mod}$
  (Theorem~\ref{thm:lines_as_modules}).  The bar functor
  $\barB^{\mathrm{ord}} \colon \cA\text{-mod}_{E_1}
  \rightleftarrows
  \cA^!\text{-comod}_{E_1} \cocolon \Omega^{\mathrm{ord}}$
  is the open-color bar-cobar adjunction.
\item \emph{Interaction.}
  The holomorphic weight filtration on $\SCchtop$ induces a
  spectral sequence $E_r$ converging to
  $H^\ast(\barB^{\SCchtop}(\cA))$ with:
  \begin{itemize}
  \item $E_0$-page: the two colors decoupled
    (Proposition~\ref{prop:gr-chiral});
  \item $d_1$-differentials encoding the classical $r$-matrix
    $r_0/z$ (leading residue of the $R$-matrix);
  \item $d_r$-differentials encoding the $r$-th OPE pole
    coefficient of $R(z)$;
  \item convergence to the full $\SCchtop$-coalgebra structure
    on~$\barB(\cA)$.
  \end{itemize}
\item \emph{Involution.}
  $(-)^{!!} \simeq \id$ at the $\SCchtop$ level:
  $\Omega^{\SCchtop}(\barB^{\SCchtop}(\cA)) \simeq \cA$
  as $\SCchtop$-algebras.  This upgrades the separate
  closed-color and open-color involutions to a single
  two-color involution.
\end{enumerate}
\end{theorem}

\begin{proof}
By Theorem~\ref{thm:homotopy-Koszul}, the operad $\SCchtop$
is homotopy-Koszul.  By Theorem~\ref{thm:SC-self-duality},
$(\SCchtop)^! \simeq \SCchtop$: the operad is Koszul
self-dual.  The general theory of homotopy-Koszul
operads (Loday--Vallette \cite{LV12}, Chapter~10; generalized
to two-colored operads by Vallette~\cite{Val07}) gives a Quillen
equivalence between $\SCchtop$-algebras and
$\SCchtop$-coalgebras via the operadic bar-cobar
adjunction (self-duality converts the
$(\SCchtop)^!$-coalgebras of the general theory into
$\SCchtop$-coalgebras).  We now project to each color.

\medskip
\noindent\textbf{(i) Closed color.}
The closed-color projection of $\barB^{\SCchtop}(\cA)$ retains
only the $\FM_k(\C)$-labeled operations, giving $\barBch(\cA)$.
The Quillen equivalence restricted to closed-color modules is
the chiral bar-cobar adjunction, which is Theorems~A and~B
of Vol~I specialized to genus~$0$ and the standard curve
$X = \C$.  The key input: the closed-color factor of
$\mathrm{gr}\;\SCchtop$ is the chiral operad
$\mathsf{P}_{\mathrm{ch}}$, which is Koszul
(Francis--Gaitsgory~\cite{FG12}).

\medskip
\noindent\textbf{(ii) Open color.}
The open-color projection retains only the
$E_1(m)$-labeled operations, giving the ordered bar complex
$\barB^{\mathrm{ord}}(\cA)$.  The Quillen equivalence
restricted to open-color modules identifies
$\cA^!$-modules with line operators, recovering
Theorem~\ref{thm:lines_as_modules}.

\medskip
\noindent\textbf{(iii) Interaction.}
The holomorphic weight filtration $F^\bullet$ on $\SCchtop$ is
the filtration by pole order on $\FM_k(\C)$.  By
Proposition~\ref{prop:gr-chiral},
$\mathrm{gr}^F\,\SCchtop \cong \mathsf{P}_{\mathrm{ch}}
\amalg E_1$ with no mixed operations.  The spectral sequence of
the filtered complex $(\barB^{\SCchtop}(\cA), F^\bullet)$
therefore has $E_0 = \barBch(\cA) \otimes
\barB^{\mathrm{ord}}(\cA)$ (decoupled).

The $d_1$ differential is the leading mixed-color correction:
it acts on bar degree~$2$ elements $[a_1 | a_2]$ by extracting
the simple pole $\Res_{z=0} \{a_{1\,\lambda}\, a_2\}/z$,
which is exactly the classical $r$-matrix $r_0 = \Omega/z$
(Proposition~\ref{prop:field-theory-r}).  At bar degree~$k$,
$d_r$ encodes the order-$r$ pole in the OPE:
$\Res_{z=0} z^{r-1} R(z)$ restricted to $k$-fold bar elements.

Convergence: the filtration is bounded below
(pole order~$\ge 0$) and exhaustive ($\FM_k(\C)$ has finite
pole orders for finite $k$), so the spectral sequence converges
by the standard bounded-filtration theorem.

\medskip
\noindent\textbf{(iv) Involution.}
The two-color bar-cobar composite
$\Omega^{\SCchtop} \circ \barB^{\SCchtop}$ is the identity on
the homotopy category of $\SCchtop$-algebras by the general
Koszul duality involution (Loday--Vallette, Theorem~10.1.13).
Since the closed-color and open-color projections of this
composite recover the separate involutions $\Omegach \circ
\barBch \simeq \id$ and
$\Omega^{\mathrm{ord}} \circ \barB^{\mathrm{ord}} \simeq
\id$, the two-color involution refines both.
\end{proof}

\subsubsection*{The spectral parameter from the closed color}

\begin{theorem}[Spectral parameter from the closed color;
\ClaimStatusProvedHere]
\label{thm:spectral-parameter-from-closed-color}
\index{spectral parameter!closed-color origin}
\index{R-matrix!two-color origin}
The spectral parameter $z$ in $r(z)$ is the coordinate on
$\FM_2(\C)$ in the closed-color factor of $\SCchtop$.
\begin{enumerate}[label=\textup{(\roman*)}]
\item The $E_1$ structure sees ordering on $\R_t$: the
  deconcatenation coproduct $\Delta$ on
  $\barB^{\mathrm{ord}}(\cA)$ is a map of graded vector spaces,
  independent of any spectral parameter.
\item The chiral structure sees holomorphic separation
  $z = z_1 - z_2$ on $\FM_2(\C)$: the bar differential
  $d_{\barB}$ on $\barBch(\cA)$ depends meromorphically on $z$.
\item The $R$-matrix lives in the \emph{mixed sector}:
  holomorphic weight $\ge 1$ (closed color contributes poles
  in $z$) and open-color degree~$2$ (acts on pairs of line
  operators via $\Delta$).
\end{enumerate}
\end{theorem}

\begin{proof}
Item~(i): the deconcatenation coproduct
$\Delta \colon T^c(V) \to T^c(V) \otimes T^c(V)$ is the
standard unshuffle coproduct on the cofree coalgebra, which is
defined purely combinatorially (Loday--Vallette, \S1.2).  No
geometric data enters; $\Delta$ is the identity of the
open-color $E_1$-coalgebra.

Item~(ii): the bar differential $d_{\barB}$ is constructed
from configuration integrals over $\FM_k(\C)$
(Sections~\ref{sec:BV_construction}--\ref{sec:FM_calculus}).
The fiber of $\FM_2(\C)$ over the diagonal is parametrized by
the tangent direction, whose coordinate is $z = z_1 - z_2$.
The OPE $a_1(z_1) a_2(z_2) \sim \sum_n C_n (z_1-z_2)^{-n-1}$
makes $d_{\barB}$ meromorphic in $z$.

Item~(iii): in the holomorphic weight filtration, an element of
$F^p$ has pole order $\ge p$ on $\FM_k(\C)$.  The identity
(pole order~$0$) is open-color-only; the $R$-matrix $r(z)$
has pole order~$1$ (the simple pole $r_0/z$) and acts on pairs
(bar degree~$2$).  Hence $r(z) \in F^1 \cap
\barB^2(\cA) \otimes \barB^2(\cA)$: mixed-sector data.
\end{proof}

\subsubsection*{The spectral sequence for the $R$-matrix}

\begin{proposition}[Spectral sequence for the $R$-matrix;
\ClaimStatusProvedHere]
\label{prop:spectral-sequence-r-matrix}
\index{spectral sequence!R-matrix}
\index{R-matrix!spectral sequence}
The holomorphic weight filtration on $\barB^{\SCchtop}(\cA)$
induces a spectral sequence $\{E_r, d_r\}_{r \ge 0}$:
\begin{enumerate}[label=\textup{(\roman*)}]
\item $E_0 = \barBch(\cA) \otimes \barB^{\mathrm{ord}}(\cA)$
  (the two colors decoupled, by
  Proposition~\ref{prop:gr-chiral}).
\item $d_1 \colon E_0^{p,q} \to E_0^{p+1,q}$ encodes the
  leading $R$-matrix coefficient: on bar degree~$2$ elements,
  $d_1([a_1 | a_2]) = r_0(a_1, a_2)/z$ where
  $r_0 = \sum_{ij} \Omega^{ij}\, e_i \otimes e_j$ is the
  Casimir (Proposition~\ref{prop:field-theory-r}).
\item At the $r$-th page, $d_r$ encodes the coefficient of
  $(z_1-z_2)^{-r-1}$ in the OPE: the $r$-th order correction
  to the $R$-matrix.
\item The spectral sequence converges:
  $E_\infty \cong H^\ast(\barB^{\SCchtop}(\cA))$.
\end{enumerate}
\end{proposition}

\begin{proof}
The spectral sequence is that of the filtered chain complex
$(\barB^{\SCchtop}(\cA), F^\bullet, d^{\SCchtop})$ where
$F^p$ consists of elements with holomorphic pole order $\ge p$.
The differential $d^{\SCchtop}$ splits as
$d^{\SCchtop} = d_0 + d_1 + d_2 + \cdots$ where $d_r$ raises
pole order by exactly $r$.  At $E_0$: $d_0$ is the sum of the
closed-color and open-color differentials, which are
independent by Proposition~\ref{prop:gr-chiral}.  The
$d_1$-differential is the leading mixed-color contribution:
from the simple pole in the OPE $a_1(z) a_2(0) \sim
C_0/(z) + \mathrm{regular}$, this is $r_0/z$.

Convergence: the filtration is bounded below ($F^0 =
\barB^{\SCchtop}(\cA)$, $F^p = 0$ for $p$ exceeding the
maximal pole order on $\FM_k(\C)$ at fixed bar degree~$k$)
and complete.  The spectral sequence of a bounded-below
exhaustive filtration on a cochain complex converges
(Weibel~\cite{Wei94}, Theorem~5.5.1).
\end{proof}

\begin{corollary}[$\Ainf$ Yang--Baxter as convergence;
\ClaimStatusProvedHere]
\label{cor:ainfty-ybe-as-convergence}
\index{Yang--Baxter equation!$A_\infty$}
\index{Yang--Baxter equation!spectral sequence convergence}
The $\Ainf$ Yang--Baxter equation
(Theorem~\ref{thm:spectral_R_YBE}) is the convergence condition
of the spectral sequence $E_r$ at bar degree~$3$: the tower of
differentials $d_1, d_2, \ldots$ acting on three-fold bar
elements $[a_1 | a_2 | a_3]$ converges if and only if the
full $R$-matrix $R(z)$ satisfies the quantum Yang--Baxter
equation on triple tensor products.
\end{corollary}

\begin{proof}
At bar degree~$3$, the $E_r$-page carries the data of three
pairwise interactions $(12)$, $(13)$, $(23)$ at holomorphic
separations $z_{12}$, $z_{13}$, $z_{23}$.  The $d_r$
differentials at this bar degree are:
\begin{align*}
d_1([a_1 | a_2 | a_3])
&= r_0^{(12)}/z_{12} + r_0^{(13)}/z_{13}
   + r_0^{(23)}/z_{23}, \\
d_r([a_1 | a_2 | a_3])
&= \text{order-$r$ pole corrections to each pair}.
\end{align*}
The condition $(d^{\SCchtop})^2 = 0$ on bar degree~$3$
elements reads, at each pole order:
\begin{itemize}
\item Order $2$: $[r_0^{(12)}, r_0^{(13)}]
  + [r_0^{(12)}, r_0^{(23)}]
  + [r_0^{(13)}, r_0^{(23)}] = 0$
  \quad (the classical YBE);
\item Order $r+1$: the correction from $d_r$ compensates
  the failure of $d_1^2 + \cdots + d_{r-1}^2$ at the previous
  stage.
\end{itemize}
The tower converges if and only if these corrections
sum to zero at all orders, which is exactly the quantum YBE
for $R(z) = \id + \hh\,r_0/z + O(z^{-2})$.
\end{proof}

\subsubsection*{The $\SCchtop$-algebra structure on $\cA^!$}

\begin{theorem}[Dual $\SCchtop$-algebra;
\ClaimStatusProvedHere]
\label{thm:dual-sc-algebra}
\index{Swiss-cheese operad!dual algebra structure}
\index{Koszul dual!$\SCchtop$-algebra}
\index{Sklyanin bracket!as closed-color cobar}
\index{Yangian!as open-color cobar}
The Koszul dual $\cA^!$ inherits a canonical
$\SCchtop$-algebra structure:
\begin{enumerate}[label=\textup{(\roman*)}]
\item \emph{Closed color:} the chiral cobar construction
  $\Omegach(\barBch(\cA))$ endows $\cA^!$ with its Sklyanin
  $\lambda$-bracket
  (equation~\eqref{eq:sklyanin-bracket}).
\item \emph{Open color:} the ordered cobar construction
  $\Omega^{\mathrm{ord}}(\barB^{\mathrm{ord}}(\cA))$ endows
  $\cA^!$ with its Yangian product.
\item \emph{Compatibility:} $d_{\barB}$ is a coderivation
  of~$\Delta$ in $\barB^{\SCchtop}(\cA)$.  Dually, the
  Sklyanin bracket and the Yangian product satisfy the
  Swiss-cheese composition law: the $\lambda$-bracket is a
  derivation of the Yangian product up to homotopy controlled
  by the mixed-color operations of $\SCchtop$.
\end{enumerate}
The $\SCchtop$-algebra on $\cA^!$ is Koszul dual to the
$\SCchtop$-coalgebra on $\barB^{\SCchtop}(\cA)$.
\end{theorem}

\begin{proof}
By the Koszul self-duality of $\SCchtop$
(Theorem~\ref{thm:SC-self-duality}), an
$\SCchtop$-coalgebra structure on $\barB(\cA)$ dualizes to an
$\SCchtop$-algebra structure on
$\cA^! = \Hom(\barB(\cA), k)$---this is the operadic
dualization of Step~2 in the proof of
Theorem~\ref{thm:Koszul_dual_Yangian}.  The recognition
theorem (Theorem~\ref{thm:yangian-recognition}) identifies
this $\SCchtop$-algebra structure as a dg-shifted Yangian;
the present theorem unpacks the two color projections and
their compatibility.
The two projections:

(i) The closed-color projection of the $\SCchtop$-algebra
structure maps are
$C_\ast(\FM_k(\C)) \otimes (\cA^!)^{\otimes k} \to \cA^!$.
At $k = 2$, the $\FM_2(\C)$-integral with meromorphic kernel
gives the $\lambda$-bracket; the Sklyanin form
\eqref{eq:sklyanin-bracket} is the explicit residue computation
(Proposition~\ref{prop:field-theory-r}).

(ii) The open-color projection is
$E_1(m) \otimes (\cA^!)^{\otimes m} \to \cA^!$: the
$E_1$-algebra structure is the Yangian product, i.e.\ the
concatenation algebra structure dual to deconcatenation.

(iii) The mixed-color composition maps of $\SCchtop$ enforce
compatibility: the bar differential $d_{\barB}$, being a
coderivation of $\Delta$, dualizes to the statement that the
$\lambda$-bracket (from $d_{\barB}$) is a derivation of the
product (from $\Delta$) up to higher $\SCchtop$-homotopies.
\end{proof}

\begin{remark}[This is not an accident]
\label{rem:not-accidental}
\index{Yangian!not accidental}
The Yangian $Y_\hh(\fg)$ is \emph{not} ``an associative algebra
that also happens to carry a Poisson vertex algebra structure.''
It is an $\SCchtop$-algebra
(Theorem~\ref{thm:yangian-recognition}): the RTT relations
encode the closed color ($\FM_2(\C)$-integrals), the Yangian
product encodes the open color ($E_1$-composition), and the
$\lambda$-brackets are the closed-color OPE\@.  The Drinfeld
presentation separates these into ``Yangian generators'' (open)
and ``spectral parameter dependence'' (closed); the
$\SCchtop$-algebra structure reunifies them.  The self-duality
$(\SCchtop)^! \simeq \SCchtop$
(Theorem~\ref{thm:SC-self-duality}) is the structural reason
the reunification works: dualizing the correspondence does not
change the type of algebra.
\end{remark}

\subsubsection*{Genus tower: asymmetry between the two colors}

\begin{theorem}[Genus tower asymmetry;
\ClaimStatusProvedHere]
\label{thm:genus-tower-asymmetry}
\index{genus tower!two-color asymmetry}
\index{curvature!closed-color only}
At genus $g \ge 1$, the two colors behave asymmetrically:
\begin{enumerate}[label=\textup{(\roman*)}]
\item \emph{Closed color is curved.}
  $\dfib^{\,2} = \kappa(\cA) \cdot \omega_g$
  (Theorem~\ref{thm:mc5-genus-one-bridge} at $g=1$;
  Vol~I, Theorem~C at all genera).  The bar complex
  $\barBch(\cA)$ on $\Sigma_g$ has non-vanishing curvature
  proportional to the modular characteristic $\kappa(\cA)$
  and the Arakelov form $\omega_g$.
\item \emph{Open color is flat.}
  The deconcatenation coproduct $\Delta$ is coassociative at
  \emph{all} genera: $(\Delta \otimes \id) \circ \Delta
  = (\id \otimes \Delta) \circ \Delta$ with no genus
  corrections.  The ordered bar complex
  $\barB^{\mathrm{ord}}(\cA)$ is a strict
  (uncurved) coalgebra at every genus.
\item \emph{Operadic origin.}
  The asymmetry is the ``no open-to-closed'' constraint of
  $\SCchtop$ (Definition~\ref{def:SC-colors}): closed-color
  operations depend on the complex structure of $\Sigma_g$
  via $\FM_k(\Sigma_g)$, which develops quasi-periodicity
  defects at $g \ge 1$.  Open-color operations
  depend only on $E_1(m)$, which is genus-independent since
  $\R$ is contractible at all genera.  The Arnold defect
  (quasi-periodic monodromy of $\FM_k(\Sigma_g)$) contributes
  only to the closed color.
\item \emph{Complementarity.}
  $\kappa(\cA) + \kappa(\cA^!) = K(\cA)$
  (Vol~I, Theorem~C), where $K(\cA)$ is a level-independent
  complementarity constant.  For the affine lineage
  ($K = 0$), the total $\SCchtop$-coalgebra
  is ``uncurved on average'': the closed-color curvatures of
  $\cA$ and $\cA^!$ cancel, and the full bulk theory
  (controlled by the pair $(\cA, \cA^!)$) sees no net
  genus-$g$ anomaly.  For DS reductions ($K \neq 0$), the
  pair retains residual curvature $K(\cA) \cdot \omega_g$.
\end{enumerate}
\end{theorem}

\begin{proof}
(i) is Theorem~\ref{thm:mc5-genus-one-bridge} at genus~$1$,
extended to all genera by the genus tower of Vol~I\@.  The bar
differential on $\Sigma_g$ acquires curvature from the period
matrix: $d_{\barB}^2 = \kappa(\cA) \cdot \omega_g$ where
$\omega_g$ is the canonical Arakelov $(1,1)$-form
(equation~\eqref{eq:curved-R-fact}).

(ii): The deconcatenation coproduct $\Delta$ on the cofree
coalgebra $T^c(V)$ is defined purely combinatorially by
$(v_1 | \cdots | v_n) \mapsto
\sum_{i=0}^n (v_1 | \cdots | v_i) \otimes
(v_{i+1} | \cdots | v_n)$.  This involves no integrals, no
propagators, and no curve geometry.  Coassociativity
$(\Delta \otimes \id) \circ \Delta
= (\id \otimes \Delta) \circ \Delta$ is an identity of graded
vector spaces, valid at every genus.

(iii): The $\SCchtop$ operad has closed-color spaces
$C_\ast(\FM_k(\C))$ and open-color spaces $E_1(m)$.  At
genus $g \ge 1$, replacing $\C$ by $\Sigma_g$: the space
$\FM_k(\Sigma_g)$ is no longer contractible (it has
$H^1(\Sigma_g)^{\oplus k} \ne 0$), and the propagator acquires
quasi-periodic monodromy from the period matrix.  This is
the Arnold defect, and it contributes curvature.  The space
$\R$, however, does not change with genus: line operators
are still ordered on the boundary half-line $\R_{\ge 0}$,
independent of $\Sigma_g$.  Hence $E_1(m)$ is genus-independent.

(iv): Complementarity $\kappa(\cA) + \kappa(\cA^!) = K(\cA)$
is Vol~I, Theorem~C, where $K(\cA)$ is the complementarity
constant.  At the $\SCchtop$ level: the closed-color
curvatures $\kappa(\cA) \cdot \omega_g$ and
$\kappa(\cA^!) \cdot \omega_g$ sum to $K(\cA) \cdot \omega_g$.
For the affine lineage ($K = 0$), they cancel when one passes to
the full bulk theory controlled by the Koszul pair
$(\cA, \cA^!)$.  For DS reductions ($K \neq 0$), the pair
retains residual curvature.
\end{proof}

\begin{remark}[Operadic explanation of the asymmetry]
\label{rem:asymmetry-operadic}
\index{genus tower!operadic explanation}
The asymmetry between colors at $g \ge 1$ is a direct
consequence of the operad structure.  The closed-color factor
$\FM_k(\Sigma_g)$ is a non-trivial fibration over
$\mathcal{M}_g$ with fibers that depend on the complex
structure; at $g = 0$, $\FM_k(\C)$ is a $K(\pi, 1)$ for the
pure braid group, and the bar differential is flat.  At
$g \ge 1$, the cohomology of $\FM_k(\Sigma_g)$ acquires new
classes from $H^1(\Sigma_g)$, and the bar differential squares
to $\kappa \cdot \omega_g \ne 0$.  The open-color factor
$E_1(m) \simeq \Sigma_m$ (a discrete set of orderings) is
independent of all continuous geometry and therefore of genus.
Curvature enters the two-color picture only through the
closed-color door.
\end{remark}

\begin{remark}[$\Theta_\cA$ and two colors]
\label{rem:theta-two-colors}
\index{shadow Postnikov tower!two-color decomposition}
\index{Maurer--Cartan element!two-color decomposition}
The universal Maurer--Cartan element
$\Theta_\cA \in \MC(\mathfrak{g}_\cA^{\mathrm{mod}})$
(from the bar-intrinsic Maurer--Cartan construction of Volume~I) encodes
\emph{both} colors simultaneously:
\begin{itemize}
\item \emph{Closed-color projection:} the shadow Postnikov
  tower $\kappa$, $C$, $Q$, \ldots are the finite-order
  projections of $\Theta_\cA$ onto the closed-color
  (chiral/modular) sector.  These are the invariants of the
  chiral bar complex $\barBch(\cA)$.
\item \emph{Open-color projection:} the $R$-matrix
  $r(z) = \Res^{\mathrm{coll}}_{0,2}(\Theta_\cA)$
  is the collision residue of $\Theta_\cA$ at arity~$2$
  in the open-color sector.  This is the invariant of the
  ordered bar complex $\barB^{\mathrm{ord}}(\cA)$.
\item \emph{Full MC equation:} $D \cdot \Theta_\cA
  + \tfrac{1}{2}[\Theta_\cA, \Theta_\cA] = 0$
  is a \emph{simultaneous} constraint on both colors.  The
  closed-color component gives $d_{\barB}^2 = 0$ (at
  genus~$0$) or $d_{\barB}^2 = \kappa \cdot \omega_g$ (at
  genus~$g \ge 1$).  The open-color component gives
  coassociativity of $\Delta$.  The mixed component gives
  the Yang--Baxter equation for $R(z)$.
\end{itemize}
The two-color Koszul duality theorem
(Theorem~\ref{thm:two-color-master}) unpacks the single MC
equation into its color components, making explicit the
information that $\Theta_\cA$ encodes simultaneously about
the chiral algebra and its line operators.
\end{remark}


%%%%%%%%%%%%%%%%%%%%%%%%%%%%%%%%%%%%%%%%%%%%%%%%%%%%%%%%%%%%%%%%%%%%%%%%%%%%%
% ELLIPTIC SPECTRAL DICHOTOMY AT GENUS 1
%%%%%%%%%%%%%%%%%%%%%%%%%%%%%%%%%%%%%%%%%%%%%%%%%%%%%%%%%%%%%%%%%%%%%%%%%%%%%

\subsection{Elliptic spectral dichotomy at genus~$1$}
\label{subsec:elliptic-spectral-dichotomy}
\index{elliptic spectral dichotomy}
\index{curvature!braiding coupling}
\index{genus 1!spectral braiding}

The genus-$1$ curvature $\dfib^{\,2} = \kappa(\cA)\cdot\omega_1$
(Theorem~\ref{thm:mc5-genus-one-bridge}) and the genus-$0$ spectral
$R$-matrix $R(z)$ (Theorem~\ref{thm:spectral_R_YBE}) originate from
different genera and appear to be independent structures.  At genus~$1$,
however, they meet: the spectral braiding must be defined on the
elliptic curve $E_\tau$, and the quasi-periodicity of the elliptic
propagator forces the two structures to interact.  The nature of this
interaction is controlled entirely by the pole structure of the
$\lambda$-bracket.

\begin{theorem}[Elliptic spectral dichotomy;
\ClaimStatusProvedHere]
\label{thm:elliptic-spectral-dichotomy}
Let\/ $\cA$ be a logarithmic\/ $\SCchtop$-algebra on $E_\tau \times \R$.
Let $\cA$ be the $A_\infty$ chiral algebra, and write the
$(-1)$-shifted $\lambda$-bracket on cohomology as
\[
\{a {}_\lambda b\}
\;=\;
c_0(a,b) + c_1(a,b)\,\lambda + c_2(a,b)\,\lambda^2 + \cdots\,,
\]
so that the genus-$0$ classical $r$-matrix
\textup{(}Proposition~\textup{\ref{prop:field-theory-r})} has the
Laurent expansion
\[
r(z)
\;=\;
\frac{c_0}{z} + \frac{c_1}{z^2} + \frac{c_2}{z^3} + \cdots
\qquad \text{\textup{(}formal, $z \to \infty$\textup{)}}.
\]
At genus~$1$, the classical $r$-matrix on $E_\tau$ is
\begin{equation}
\label{eq:elliptic-r-matrix}
r^{E_\tau}(z)
\;=\;
c_0 \cdot \zeta(z|\tau)
\;+\;
c_1 \cdot \wp(z|\tau)
\;+\;
\sum_{n \ge 2}
  \frac{(-1)^{n-1}}{(n-1)!}\,
  c_n \cdot \wp^{(n-2)}(z|\tau),
\end{equation}
obtained by replacing the rational functions $1/z^{n+1}$ in the
Laurent expansion of $r(z)$ by their elliptic counterparts on
$E_\tau$: the unique functions matching $1/z^{n+1} + O(z)$
near $z = 0$ and having prescribed periodicity.  Here
$\zeta(z|\tau)$, $\wp(z|\tau)$ are the Weierstrass zeta and
$\wp$-functions with periods $1, \tau$.

Then:
\begin{enumerate}[label=\textup{(\alph*)}]
\item \label{item:dichotomy-monodromy}
\textup{(Quasi-periodicity.)}
The $B$-cycle monodromy of $r^{E_\tau}$ is
\begin{equation}
\label{eq:r-monodromy}
r^{E_\tau}(z + \tau) - r^{E_\tau}(z)
\;=\;
2\eta_\tau \cdot c_0,
\end{equation}
where $\eta_\tau = \zeta(\tau/2|\tau)$ is the quasi-period.  The
$A$-cycle monodromy is $2\eta_1 \cdot c_0$.

\item \label{item:dichotomy-decoupling}
\textup{(Decoupling: Cartan type.)}
If $c_0 = 0$---equivalently, if
$\{a {}_0 b\} = 0$ for all
$a, b \in H^\bullet(\cA, Q)$---then $r^{E_\tau}(z)$ is a
meromorphic function on $E_\tau$ \textup{(}doubly
periodic\textup{)}.
The spectral braiding is independent of the $B$-cycle monodromy
that sources the curvature $\kappa(\cA) \cdot \omega_1$.  The
genus-$1$ Swiss-cheese structure splits:
\[
\text{curved}\;\mathsf{SC}^{\mathrm{ch,top},(1)}
\;\simeq\;
(\text{curved closed color})
\;\times\;
(\text{flat open color with meromorphic braiding}).
\]

\item \label{item:dichotomy-entanglement}
\textup{(Entanglement: Yangian type.)}
If $c_0 \ne 0$, then $r^{E_\tau}(z)$ is quasi-periodic with
non-trivial $B$-cycle monodromy $2\eta_\tau \cdot c_0$.  The
spectral braiding is a section of a flat connection on the universal
cover of $E_\tau$ with holonomy~$c_0$, and the genus-$1$
Swiss-cheese structure does \emph{not} split.
The braiding and the curvature share a common geometric source:
the $B$-cycle monodromy of the Weierstrass zeta function.
\end{enumerate}
\end{theorem}

\begin{proof}
\textbf{Step~1} (Elliptic propagator in the exchange diagram).
The spectral $R$-matrix $R(z)$ at genus~$0$ is computed from
a single bulk-to-boundary propagator exchange
(Proposition~\ref{prop:field-theory-r}, Step~1), with
holomorphic propagator $K^{\C}(z) = dz/z$.  At genus~$1$,
the theory on $E_\tau \times \R_{\ge 0}$ uses the elliptic
propagator
$K^{E_\tau}(z) = \zeta(z|\tau)\,dz$
(Theorem~\ref{thm:mc5-genus-one-bridge}, Step~1), which has
the same local singularity
$\zeta(z|\tau) = 1/z + O(z)$ near $z = 0$
but acquires the quasi-periodicity
\[
\zeta(z + 1|\tau) = \zeta(z|\tau) + 2\eta_1,
\qquad
\zeta(z + \tau|\tau) = \zeta(z|\tau) + 2\eta_\tau.
\]

\textbf{Step~2} (Laplace formula at genus~$1$).
The $\lambda$-bracket on cohomology is a \emph{local} object:
it depends only on the residue of the propagator at the
collision divisor $z_1 = z_2$
(Section~\ref{sec:Ainfty-to-PVA}), which is the same at
genus~$0$ and genus~$1$
(Theorem~\ref{thm:mc5-genus-one-bridge}, Step~1).  Therefore
$\{a {}_\lambda b\}^{E_\tau} = \{a {}_\lambda b\}^{\C}$,
and the coefficients $c_n$ are genus-independent.

However, the classical $r$-matrix is a \emph{global} object:
it depends on the propagator at \emph{separated} points
$z_1, z_2$.  At genus~$1$, the Laplace formula of
Proposition~\ref{prop:field-theory-r} uses the elliptic
propagator $\zeta(z|\tau)$ in place of $1/z$.  The formal
substitution $1/z^{n+1} \leadsto f_{n+1}(z|\tau)$ is
determined by matching poles:
\begin{align*}
\frac{1}{z}
  &\;\leadsto\;
  \zeta(z|\tau)
  &&\text{(quasi-periodic)}, \\
\frac{1}{z^2}
  &\;\leadsto\;
  \wp(z|\tau) = -\zeta'(z|\tau)
  &&\text{(periodic)}, \\
\frac{1}{z^{n+1}}
  &\;\leadsto\;
  \frac{(-1)^{n}}{n!}\,\wp^{(n-1)}(z|\tau)
  &&\text{(periodic, $n \ge 2$)}.
\end{align*}
These are the unique functions with pole $1/z^{n+1} + O(z)$ at
$z = 0$ that are doubly periodic for $n \ge 1$, and
quasi-periodic with prescribed quasi-periods for $n = 0$.
Substituting into $r(z) = \sum_{n \ge 0} c_n / z^{n+1}$
gives equation~\eqref{eq:elliptic-r-matrix}.

\textbf{Step~3} (Quasi-periodicity).
The monodromy~\eqref{eq:r-monodromy} follows directly:
the only non-periodic function in the
expansion~\eqref{eq:elliptic-r-matrix} is
$c_0 \cdot \zeta(z|\tau)$, since all terms involving
$\wp^{(n-1)}$ for $n \ge 1$ are doubly periodic.
The quasi-periodicity
$\zeta(z + \tau|\tau) - \zeta(z|\tau) = 2\eta_\tau$
gives~\ref{item:dichotomy-monodromy}.

\textbf{Step~4} (Common geometric source).
The same $B$-cycle quasi-period $2\eta_\tau$ of the Weierstrass
zeta function produces two distinct algebraic structures:
\begin{itemize}
\item \emph{Curvature.}
In the proof of Theorem~\ref{thm:mc5-genus-one-bridge}
(Step~3), the quasi-periodicity of $\zeta$ breaks the Arnold
relation on $\FM_3(E_\tau)$, producing the defect
$E_2(\tau) \cdot (dz_1 - dz_2) \wedge (dz_2 - dz_3)$.
The Eisenstein series $E_2(\tau)$ is related to the
quasi-period by $\eta_1 = \pi^2 E_2(\tau) / 6$.
Contracting with the OPE data extracts the
\emph{highest-pole coefficient} of the $\lambda$-bracket
(the term $c_1 \cdot \lambda$ for affine algebras), yielding
the scalar curvature
$\dfib^{\,2} = \kappa(\cA) \cdot \omega_1$.

\item \emph{Braiding monodromy.}
In the elliptic $r$-matrix~\eqref{eq:elliptic-r-matrix},
the same quasi-periodicity produces the monodromy
$\delta r = 2\eta_\tau \cdot c_0$, which extracts the
\emph{lowest-pole coefficient} $c_0 = \{a {}_0 b\}$ of the
$\lambda$-bracket.
\end{itemize}
The curvature is a \emph{scalar} (the trace of the
highest-pole coefficient), invisible to the tensor structure
of the $r$-matrix.
The braiding monodromy is a \emph{tensor} (the full
lowest-pole coefficient), carrying the Lie-algebraic structure
constants.  When $c_0 = 0$, the braiding monodromy vanishes
and the only effect of the $B$-cycle is the scalar curvature:
decoupling~\ref{item:dichotomy-decoupling}.  When
$c_0 \ne 0$, both scalar and tensorial effects are sourced
by the same geometric defect:
entanglement~\ref{item:dichotomy-entanglement}.
\end{proof}

\begin{remark}[The Coisson bracket as discriminant]
\label{rem:Coisson-discriminant}
\index{Coisson bracket!curvature-braiding discriminant}
The coefficient $c_0 = \{a {}_0 b\}$ is the $\lambda = 0$
specialization of the $\lambda$-bracket---the Coisson bracket
of Theorem~\ref{thm:pva-coisson-bridge}.  The elliptic spectral
dichotomy is therefore:
\[
\text{Coisson bracket vanishes}
\;\;\Longleftrightarrow\;\;
\text{curvature-braiding decoupling at genus~$1$}.
\]
The Coisson structure is therefore the obstruction to the independence of holomorphic
and topological deformations at genus~$1$.
\end{remark}

\begin{remark}[Verification in the standard landscape]
\label{rem:elliptic-dichotomy-examples}
The dichotomy is verified in the standard examples:
\begin{enumerate}[label=(\roman*)]
\item \emph{Heisenberg $\mathcal{H}_k$:}
  $\{J {}_\lambda J\} = k\lambda$, so $c_0 = 0$.
  The elliptic $r$-matrix is
  $r^{E_\tau}(z) = k \cdot \wp(z|\tau)$,
  doubly periodic on $E_\tau$.  Curvature
  $\dfib^{\,2} = k \cdot \omega_1$
  (Theorem~\ref{thm:rosetta-curved}) and braiding decouple.

\item \emph{Affine Kac--Moody $\widehat{\fg}_k$:}
  $\{J^a {}_\lambda J^b\} = f^{ab}_c\, J^c +
  k\,\kappa^{ab}\,\lambda$,
  so $c_0^{ab} = f^{ab}_c\, J^c \ne 0$.
  The elliptic $r$-matrix
  $r^{E_\tau}(z) = \Omega \cdot \zeta(z|\tau) +
  k\kappa \cdot \wp(z|\tau)$
  has $B$-cycle monodromy $2\eta_\tau \cdot \Omega$
  (where $\Omega = \sum_a t^a \otimes t_a$ is the quadratic
  Casimir).  Curvature and braiding entangle.

\item \emph{Virasoro $\mathrm{Vir}_c$:}
  $\{L {}_\lambda L\} = (\partial + 2\lambda)\,L +
  \tfrac{c}{12}\,\lambda^3$,
  so $c_0 = \partial L \ne 0$.
  Curvature and braiding entangle.
\end{enumerate}
\end{remark}

\begin{remark}[Connection to the elliptic KZ equation]
\label{rem:elliptic-KZ}
For $\widehat{\fg}_k$, the elliptic $r$-matrix
$r^{E_\tau}(z) = \Omega \cdot \zeta(z|\tau) + \cdots$
is the classical limit of the elliptic Knizhnik--Zamolodchikov
connection of Bernard \cite{Ber88}:
\[
\nabla_i^{\mathrm{KZ}}
\;=\;
\partial_{z_i}
\;-\;
\frac{1}{k + h^\vee}
\sum_{j \ne i}
  \Omega_{ij} \cdot \wp(z_i - z_j|\tau).
\]
The quasi-periodicity of $\zeta$ in the $r$-matrix
corresponds to the non-trivial $B$-cycle monodromy of
the KZ connection.  The curvature-braiding entanglement
at genus~$1$ is, from this perspective, the statement
that the elliptic KZ connection on $E_\tau^n$ and the
modular differential equation
$\partial_\tau \Phi = \Delta_\tau \Phi$ (governing
genus-$1$ partition functions) share the same Casimir
data~$\Omega$.  For abelian algebras ($\Omega$ scalar),
the KZ connection is flat and trivializable;
for non-abelian algebras, the holonomy representation
of $\pi_1(E_\tau)$ is irreducible.
\end{remark}

\begin{remark}[Operadic splitting obstruction]
\label{rem:operadic-splitting}
In the language of the Swiss-cheese operad, the dichotomy
states: at genus~$1$, the curved $\SCchtop$-algebra structure
on $\barB(\cA)$ can be presented as a product of a curved
(closed-color) Beilinson--Drinfeld algebra and a flat (open-color)
$E_1$-algebra if and only if $c_0 = 0$.  When $c_0 \ne 0$,
the two colors are coupled by the non-trivial $B$-cycle
holonomy: the open-color braiding ``sees'' the closed-color
monodromy through the shared quasi-periodicity of the
propagator.  The obstruction to splitting is the element
$c_0 \in H^\bullet(\cA,Q)^{\otimes 2}$, the infinitesimal
generator of the holonomy representation.
\end{remark}
